% Chapter 11. All quoted numbers come from chapters/ch11-numbers.tex, written by code/ch11/run_all.py.
% Generated by code/ch11/run_all.py. Do not edit by hand.
\newcommand{\ckFullEight}{1.65}
\newcommand{\ckFullHalf}{1.74}
\newcommand{\ckFullOne}{1.63}
\newcommand{\ckShareZero}{11.3}
\newcommand{\ckZeroMid}{11.5}
\newcommand{\ctAPeak}{4.6}
\newcommand{\ctClaw}{1.6}
\newcommand{\ctFinal}{3.0}
\newcommand{\ctPOne}{25}
\newcommand{\ctPThree}{2}
\newcommand{\ctPTwo}{18}
\newcommand{\ctPvGap}{0.593}
\newcommand{\ctR}{8}
\newcommand{\ctSOne}{2}
\newcommand{\ctSThree}{6}
\newcommand{\ctSTwo}{4}
\newcommand{\cvCommitted}{163.0}
\newcommand{\cvDpiFifteen}{1.75}
\newcommand{\cvDpiNineteen}{0.50}
\newcommand{\cvDpiSeven}{1.54}
\newcommand{\cvDpiTwentyThree}{0.05}
\newcommand{\cvDurMed}{4.1}
\newcommand{\cvDurMid}{4.2}
\newcommand{\cvDurN}{210}
\newcommand{\cvDurNMid}{44}
\newcommand{\cvDurNPre}{50}
\newcommand{\cvDurOver}{0}
\newcommand{\cvDurPre}{5.6}
\newcommand{\cvDurRatioMax}{82}
\newcommand{\cvDurShareMax}{42}
\newcommand{\cvFirstVintage}{1998}
\newcommand{\cvFunds}{485}
\newcommand{\cvNFifteen}{12}
\newcommand{\cvNNineteen}{22}
\newcommand{\cvNSeven}{23}
\newcommand{\cvNTwentyThree}{71}
\newcommand{\cvNav}{119.6}
\newcommand{\cvNm}{261}
\newcommand{\cvNmFrom}{2022}
\newcommand{\cvOldN}{77}
\newcommand{\cvOldRvpi}{3}
\newcommand{\cvOldTvpi}{1.65}
\newcommand{\cvTvpiFifteen}{2.17}
\newcommand{\cvTvpiNineteen}{1.52}
\newcommand{\cvTvpiSeven}{1.56}
\newcommand{\cvTvpiTwentyThree}{1.53}
\newcommand{\dmChSixtyFour}{32.5}
\newcommand{\dmChSixtyOne}{20.7}
\newcommand{\dmChThirtyFour}{13.5}
\newcommand{\dmChThirtyOne}{9.4}
\newcommand{\dmChaffe}{19.3}
\newcommand{\dmChaffeZero}{25.0}
\newcommand{\dmComb}{35.5}
\newcommand{\dmCp}{25}
\newcommand{\dmDloc}{20}
\newcommand{\dmLong}{61.8}
\newcommand{\dmLongOverPut}{3.2}
\newcommand{\dmLongSixFour}{137.4}
\newcommand{\dmLsSixtyFour}{137.4}
\newcommand{\dmLsSixtyOne}{57.6}
\newcommand{\dmLsThirtyFour}{57.6}
\newcommand{\dmLsThirtyOne}{26.3}
\newcommand{\dmMcMaxZ}{1.9}
\newcommand{\dmMcPaths}{200,000}
\newcommand{\dmMcSixFour}{137.4}
\newcommand{\dmMcSixFourSE}{0.5}
\newcommand{\dmMcSteps}{48}
\newcommand{\dmSig}{45}
\newcommand{\dmStake}{11.53}
\newcommand{\dmStakeCtl}{17.86}
\newcommand{\dmStakeShare}{10}
\newcommand{\dmSum}{39.3}
\newcommand{\dmT}{2}
\newcommand{\drFiveSix}{5.5}
\newcommand{\drFiveTwelve}{19.3}
\newcommand{\drFiveTwenty}{24.8}
\newcommand{\drZero}{12.5}
\newcommand{\duEx}{0.82}
\newcommand{\duTen}{7.3}
\newcommand{\jcAlpha}{3}
\newcommand{\jcBeta}{1.2}
\newcommand{\jcCross}{3.25}
\newcommand{\jcDealsPer}{12}
\newcommand{\jcDpiEnd}{1.23}
\newcommand{\jcDpiFour}{0.25}
\newcommand{\jcDpiOne}{0.00}
\newcommand{\jcDpiTwo}{0.00}
\newcommand{\jcFunds}{4,000}
\newcommand{\jcIrrEnd}{5.2}
\newcommand{\jcIrrFour}{2.0}
\newcommand{\jcIrrHiTwo}{19.3}
\newcommand{\jcIrrLoTwo}{\ensuremath{-}29.9}
\newcommand{\jcIrrOne}{\ensuremath{-}22.3}
\newcommand{\jcIrrTwo}{\ensuremath{-}6.3}
\newcommand{\jcLife}{12}
\newcommand{\jcMuLessVar}{4.44}
\newcommand{\jcMuM}{7}
\newcommand{\jcRankEight}{0.94}
\newcommand{\jcRankFive}{0.72}
\newcommand{\jcRankThree}{0.40}
\newcommand{\jcSigE}{35}
\newcommand{\jcSigM}{16}
\newcommand{\jcTopEight}{86}
\newcommand{\jcTopFive}{68}
\newcommand{\jcTopThree}{48}
\newcommand{\jcTvpiEnd}{1.23}
\newcommand{\jcTvpiFour}{1.04}
\newcommand{\jcTvpiOne}{0.90}
\newcommand{\jcTvpiTwo}{0.94}
\newcommand{\lxBrD}{126.8}
\newcommand{\lxBrDs}{35}
\newcommand{\lxBrF}{10.0}
\newcommand{\lxBrG}{213.8}
\newcommand{\lxBrGs}{59}
\newcommand{\lxBrM}{30.2}
\newcommand{\lxBrMs}{8}
\newcommand{\lxCapex}{3.5}
\newcommand{\lxDa}{3}
\newcommand{\lxDz}{250}
\newcommand{\lxEqt}{620.8}
\newcommand{\lxEqz}{260}
\newcommand{\lxEt}{70.9}
\newcommand{\lxEvt}{744.0}
\newcommand{\lxEvz}{500}
\newcommand{\lxEz}{50}
\newcommand{\lxFees}{2}
\newcommand{\lxGain}{360.8}
\newcommand{\lxGfirst}{208.6}
\newcommand{\lxGridHiHi}{15.0}
\newcommand{\lxGridHiLo}{4.2}
\newcommand{\lxGridLoHi}{35.7}
\newcommand{\lxGridLoLo}{22.9}
\newcommand{\lxGridMax}{12}
\newcommand{\lxGridMin}{8}
\newcommand{\lxI}{8}
\newcommand{\lxIrr}{19.0}
\newcommand{\lxIrrU}{13.4}
\newcommand{\lxK}{1.338}
\newcommand{\lxLamB}{49.0}
\newcommand{\lxLamBase}{50}
\newcommand{\lxLamEx}{60}
\newcommand{\lxLev}{5.0}
\newcommand{\lxLevExit}{1.74}
\newcommand{\lxMHi}{2.99}
\newcommand{\lxMLo}{0.74}
\newcommand{\lxMcBeat}{86.8}
\newcommand{\lxMcLossL}{3.4}
\newcommand{\lxMcLossU}{0.5}
\newcommand{\lxMcMedL}{15.1}
\newcommand{\lxMcMedU}{11.0}
\newcommand{\lxMcN}{20,000}
\newcommand{\lxMcQL}{1.5}
\newcommand{\lxMcQU}{3.8}
\newcommand{\lxMcSdL}{8.0}
\newcommand{\lxMcSdU}{4.6}
\newcommand{\lxMcUAbove}{86.8}
\newcommand{\lxMcWipe}{0.00}
\newcommand{\lxMfirst}{219.0}
\newcommand{\lxMoic}{2.39}
\newcommand{\lxMoicU}{1.87}
\newcommand{\lxMt}{10.5}
\newcommand{\lxMz}{10.0}
\newcommand{\lxNdt}{123.2}
\newcommand{\lxNwc}{10}
\newcommand{\lxRate}{8}
\newcommand{\lxRev}{250}
\newcommand{\lxT}{5}
\newcommand{\lxTCapexFive}{11.3}
\newcommand{\lxTCapexFour}{10.8}
\newcommand{\lxTCapexOne}{9.3}
\newcommand{\lxTCapexThree}{10.3}
\newcommand{\lxTCapexTwo}{9.8}
\newcommand{\lxTDebtFive}{123.2}
\newcommand{\lxTDebtFour}{156.9}
\newcommand{\lxTDebtOne}{233.0}
\newcommand{\lxTDebtThree}{186.8}
\newcommand{\lxTDebtTwo}{212.1}
\newcommand{\lxTDnwcFive}{1.2}
\newcommand{\lxTDnwcFour}{1.5}
\newcommand{\lxTDnwcOne}{1.5}
\newcommand{\lxTDnwcThree}{1.4}
\newcommand{\lxTDnwcTwo}{1.6}
\newcommand{\lxTEbitdaFive}{70.9}
\newcommand{\lxTEbitdaFour}{68.1}
\newcommand{\lxTEbitdaOne}{54.3}
\newcommand{\lxTEbitdaThree}{63.4}
\newcommand{\lxTEbitdaTwo}{59.0}
\newcommand{\lxTFcfFive}{33.6}
\newcommand{\lxTFcfFour}{29.9}
\newcommand{\lxTFcfOne}{17.0}
\newcommand{\lxTFcfThree}{25.3}
\newcommand{\lxTFcfTwo}{20.9}
\newcommand{\lxTInterestFive}{12.6}
\newcommand{\lxTInterestFour}{14.9}
\newcommand{\lxTInterestOne}{20.0}
\newcommand{\lxTInterestThree}{17.0}
\newcommand{\lxTInterestTwo}{18.6}
\newcommand{\lxTRevFive}{322.1}
\newcommand{\lxTRevFour}{309.7}
\newcommand{\lxTRevOne}{265.0}
\newcommand{\lxTRevThree}{294.9}
\newcommand{\lxTRevTwo}{280.9}
\newcommand{\lxTTaxFive}{12.2}
\newcommand{\lxTTaxFour}{11.0}
\newcommand{\lxTTaxOne}{6.6}
\newcommand{\lxTTaxThree}{9.4}
\newcommand{\lxTTaxTwo}{8.0}
\newcommand{\lxTau}{25}
\newcommand{\lxUHi}{2.0}
\newcommand{\lxULo}{1.1}
\newcommand{\lxUb}{1.87}
\newcommand{\mgLbFortyFiveFour}{0.79}
\newcommand{\mgLbFortyFiveOne}{8.39}
\newcommand{\mgLbFortyFiveTwo}{5.46}
\newcommand{\mgLbSixtyFour}{n.a.}
\newcommand{\mgLbSixtyOne}{6.06}
\newcommand{\mgLbSixtyTwo}{1.76}
\newcommand{\mgLbThirtyFour}{6.06}
\newcommand{\mgLbThirtyOne}{10.53}
\newcommand{\mgLbThirtyTwo}{8.77}
\newcommand{\mgMissing}{1}
\newcommand{\mgPutFortyFiveFour}{10.97}
\newcommand{\mgPutFortyFiveOne}{12.14}
\newcommand{\mgPutFortyFiveTwo}{11.53}
\newcommand{\mgPutSixtyFour}{9.64}
\newcommand{\mgPutSixtyOne}{11.34}
\newcommand{\mgPutSixtyTwo}{10.47}
\newcommand{\mgPutThirtyFour}{12.35}
\newcommand{\mgPutThirtyOne}{12.95}
\newcommand{\mgPutThirtyTwo}{12.62}
\newcommand{\mkAbove}{69}
\newcommand{\mkGstar}{4.06}
\newcommand{\mkImpMed}{17.8}
\newcommand{\mkN}{81}
\newcommand{\mkObsMed}{25.0}
\newcommand{\mkRatioMed}{1.34}
\newcommand{\mkRho}{\ensuremath{-}0.62}
\newcommand{\mkRhoG}{0.27}
\newcommand{\mkRhoTwo}{\ensuremath{-}0.35}
\newcommand{\mkRhoW}{0.51}
\newcommand{\mkRocMed}{15.6}
\newcommand{\mkWmed}{7.13}
\newcommand{\nmChange}{10}
\newcommand{\nmGain}{\ensuremath{-}0.4}
\newcommand{\nmLegal}{0.9}
\newcommand{\nmMargin}{16.8}
\newcommand{\nmNorm}{20.2}
\newcommand{\nmPay}{1.2}
\newcommand{\nmRent}{\ensuremath{-}0.5}
\newcommand{\nmRep}{18.4}
\newcommand{\nmRev}{120}
\newcommand{\nmStock}{0.6}
\newcommand{\pkAbove}{40}
\newcommand{\pkAboveBetaOne}{35}
\newcommand{\pkBetaBoost}{0.89}
\newcommand{\pkBetaHalf}{1.014}
\newcommand{\pkBetaOneHalf}{0.986}
\newcommand{\pkBetaSimTwo}{0.974}
\newcommand{\pkBetaSimTwoSE}{0.001}
\newcommand{\pkBetaTwo}{0.973}
\newcommand{\pkBetaTwoCont}{0.976}
\newcommand{\pkBetaTwoErp}{1.087}
\newcommand{\pkCorr}{0.990}
\newcommand{\pkDaMed}{\ensuremath{-}1.3}
\newcommand{\pkErp}{4.23}
\newcommand{\pkFakeAlpha}{\ensuremath{-}0.275}
\newcommand{\pkMed}{0.95}
\newcommand{\pkMedBetaOne}{0.92}
\newcommand{\pkMu}{7}
\newcommand{\pkPrem}{2.01}
\newcommand{\pkRfCont}{4.93}
\newcommand{\pkSig}{16}
\newcommand{\pkSnipAlpha}{13.09}
\newcommand{\pkSnipIrr}{20.47}
\newcommand{\pkSnipPme}{1.6624}
\newcommand{\pkSnipTvpi}{2.10}
\newcommand{\pkT}{5}
\newcommand{\pkVar}{2.56}
\newcommand{\trBottomBottomFive}{62}
\newcommand{\trBottomBottomThree}{41}
\newcommand{\trBottomTopFive}{1}
\newcommand{\trBottomTopThree}{7}
\newcommand{\trMissFive}{0}
\newcommand{\trMissThree}{0}
\newcommand{\trTopBottomFive}{3}
\newcommand{\trTopBottomThree}{12}
\newcommand{\trTopTopFive}{68}
\newcommand{\trTopTopThree}{48}
\newcommand{\vdCapexFive}{5.16}
\newcommand{\vdCapexFour}{5.01}
\newcommand{\vdCapexOne}{4.41}
\newcommand{\vdCapexThree}{4.82}
\newcommand{\vdCapexTwo}{4.63}
\newcommand{\vdDaFive}{4.42}
\newcommand{\vdDaFour}{4.29}
\newcommand{\vdDaOne}{3.78}
\newcommand{\vdDaThree}{4.13}
\newcommand{\vdDaTwo}{3.97}
\newcommand{\vdDfFive}{0.680}
\newcommand{\vdDfFour}{0.734}
\newcommand{\vdDfOne}{0.926}
\newcommand{\vdDfThree}{0.793}
\newcommand{\vdDfTwo}{0.857}
\newcommand{\vdDwcFive}{0.52}
\newcommand{\vdDwcFour}{0.66}
\newcommand{\vdDwcOne}{0.72}
\newcommand{\vdDwcThree}{0.64}
\newcommand{\vdDwcTwo}{0.76}
\newcommand{\vdEbitFive}{20.39}
\newcommand{\vdEbitFour}{19.79}
\newcommand{\vdEbitOne}{17.43}
\newcommand{\vdEbitThree}{19.03}
\newcommand{\vdEbitTwo}{18.30}
\newcommand{\vdEbitdaFive}{24.81}
\newcommand{\vdEbitdaFour}{24.09}
\newcommand{\vdEbitdaOne}{21.21}
\newcommand{\vdEbitdaThree}{23.16}
\newcommand{\vdEbitdaTwo}{22.27}
\newcommand{\vdFcfFive}{14.04}
\newcommand{\vdFcfFour}{13.47}
\newcommand{\vdFcfOne}{11.72}
\newcommand{\vdFcfThree}{12.95}
\newcommand{\vdFcfTwo}{12.31}
\newcommand{\vdNopatFive}{15.29}
\newcommand{\vdNopatFour}{14.85}
\newcommand{\vdNopatOne}{13.07}
\newcommand{\vdNopatThree}{14.28}
\newcommand{\vdNopatTwo}{13.73}
\newcommand{\vdPvFive}{9.54}
\newcommand{\vdPvFour}{9.89}
\newcommand{\vdPvOne}{10.85}
\newcommand{\vdPvSum}{51.10}
\newcommand{\vdPvThree}{10.27}
\newcommand{\vdPvTv}{152.50}
\newcommand{\vdPvTwo}{10.55}
\newcommand{\vdRevFive}{147.39}
\newcommand{\vdRevFour}{143.10}
\newcommand{\vdRevOne}{126.00}
\newcommand{\vdRevThree}{137.59}
\newcommand{\vdRevTwo}{132.30}
\newcommand{\vlCapEarn}{212.0}
\newcommand{\vlEq}{178.6}
\newcommand{\vlEv}{203.6}
\newcommand{\vlEvEbitda}{10.08}
\newcommand{\vlFcfFive}{14.04}
\newcommand{\vlFcfFour}{13.47}
\newcommand{\vlFcfOne}{11.72}
\newcommand{\vlFcfThree}{12.95}
\newcommand{\vlFcfTwo}{12.31}
\newcommand{\vlG}{2.5}
\newcommand{\vlIndMult}{16.22}
\newcommand{\vlMktEv}{327.6}
\newcommand{\vlMktGap}{61}
\newcommand{\vlNetDebt}{25}
\newcommand{\vlNopatSix}{15.67}
\newcommand{\vlReinv}{20.8}
\newcommand{\vlRoc}{12}
\newcommand{\vlShield}{16.9}
\newcommand{\vlShieldShare}{8.3}
\newcommand{\vlTauD}{15.3}
\newcommand{\vlToCapexDn}{\ensuremath{-}1}
\newcommand{\vlToCapexUp}{1}
\newcommand{\vlToGrowthDn}{\ensuremath{-}2.9}
\newcommand{\vlToGrowthUp}{3.5}
\newcommand{\vlToMarginDn}{\ensuremath{-}7}
\newcommand{\vlToMarginUp}{7}
\newcommand{\vlToNearDn}{\ensuremath{-}8}
\newcommand{\vlToNearUp}{8}
\newcommand{\vlToRocDn}{\ensuremath{-}7}
\newcommand{\vlToRocUp}{4}
\newcommand{\vlToWaccDn}{\ensuremath{-}15}
\newcommand{\vlToWaccUp}{22}
\newcommand{\vlTv}{224.4}
\newcommand{\vlTvNopat}{14.32}
\newcommand{\vlTvNopatAtW}{12.45}
\newcommand{\vlTvNopatHiRoc}{18.08}
\newcommand{\vlTvShare}{75}
\newcommand{\vlVu}{186.7}
\newcommand{\wfCarry}{20}
\newcommand{\wfCatch}{6.57}
\newcommand{\wfClaw}{0.00}
\newcommand{\wfCommit}{100}
\newcommand{\wfCost}{10}
\newcommand{\wfDeals}{8}
\newcommand{\wfDpiPaid}{139.9}
\newcommand{\wfFeeInv}{2}
\newcommand{\wfFeePost}{1.5}
\newcommand{\wfFees}{11.7}
\newcommand{\wfFirstCarry}{8}
\newcommand{\wfGiHigh}{27.1}
\newcommand{\wfGiLow}{4.6}
\newcommand{\wfGiMid}{16.2}
\newcommand{\wfGmHigh}{2.5}
\newcommand{\wfGmLow}{1.2}
\newcommand{\wfGmMid}{1.8}
\newcommand{\wfGpShare}{20.00}
\newcommand{\wfGpSplit}{5.50}
\newcommand{\wfGpTotal}{12.07}
\newcommand{\wfGross}{152.0}
\newcommand{\wfGrossAtHurdle}{11.8}
\newcommand{\wfGrossIrr}{17.9}
\newcommand{\wfGrossMoic}{1.90}
\newcommand{\wfH}{8}
\newcommand{\wfInvYears}{5}
\newcommand{\wfLpSplit}{21.99}
\newcommand{\wfLpTotal}{139.9}
\newcommand{\wfMoicAtHurdle}{1.56}
\newcommand{\wfNetIrr}{12.0}
\newcommand{\wfNetIrrFee}{13.9}
\newcommand{\wfNetTvpi}{1.53}
\newcommand{\wfNiHigh}{19.8}
\newcommand{\wfNiLow}{1.2}
\newcommand{\wfNiMid}{10.6}
\newcommand{\wfPaidIn}{91.7}
\newcommand{\wfPref}{26.29}
\newcommand{\wfProfit}{60.3}
\newcommand{\wfRoc}{91.7}
\newcommand{\wfSpreadHigh}{7.3}
\newcommand{\wfSpreadLow}{3.5}
\newcommand{\wfSpreadMid}{5.6}
\newcommand{\wfTmCap}{120}
\newcommand{\wfTmClawFreq}{66}
\newcommand{\wfTmClawMed}{45}
\newcommand{\wfTmDeals}{12}
\newcommand{\wfTmExMean}{1.35}
\newcommand{\wfTmExMed}{1.30}
\newcommand{\wfTmExNine}{2.18}
\newcommand{\wfTmFunds}{20,000}
\newcommand{\wfTmMu}{10}
\newcommand{\wfTmNoCarry}{14.8}
\newcommand{\wfTmR}{8}
\newcommand{\wfTmRatio}{27.5}
\newcommand{\wfTmSig}{45}
\newcommand{\wkBetaL}{0.9636}
\newcommand{\wkBrHam}{1.147}
\newcommand{\wkBrHp}{1.200}
\newcommand{\wkBuHam}{0.868}
\newcommand{\wkBuHp}{0.840}
\newcommand{\wkDvInd}{12.8}
\newcommand{\wkErp}{4.23}
\newcommand{\wkErpFirst}{1961}
\newcommand{\wkErpMax}{6.45}
\newcommand{\wkErpMaxYear}{1979}
\newcommand{\wkErpMean}{4.25}
\newcommand{\wkErpMin}{2.05}
\newcommand{\wkErpMinYear}{1999}
\newcommand{\wkFirms}{105}
\newcommand{\wkRd}{6.74}
\newcommand{\wkReHam}{9.84}
\newcommand{\wkReHp}{10.07}
\newcommand{\wkReMe}{9.30}
\newcommand{\wkRf}{4.99}
\newcommand{\wkRu}{8.54}
\newcommand{\wkSpread}{1.75}
\newcommand{\wkTarget}{30}
\newcommand{\wkTax}{25}
\newcommand{\wkTbond}{4.18}
\newcommand{\wkWHam}{8.41}
\newcommand{\wkWHp}{8.56}
\newcommand{\wkWme}{8.03}
\newcommand{\wlGp}{18.00}
\newcommand{\wlLp}{172.00}
\newcommand{\wlProfit}{90}
\newcommand{\wrHamSixty}{8.15}
\newcommand{\wrHamZero}{8.66}
\newcommand{\wrHpSixty}{8.58}
\newcommand{\wrHpZero}{8.54}
\newcommand{\wrMeSixty}{7.52}
\newcommand{\wrMeZero}{8.54}
\newcommand{\wrSpreadSixty}{1.07}
\newcommand{\wrSpreadTarget}{0.53}
\newcommand{\yaBalSix}{83.10}
\newcommand{\yaBalThree}{65.97}
\newcommand{\yaCatch}{10.78}
\newcommand{\yaGp}{16.00}
\newcommand{\yaLp}{164.00}
\newcommand{\yaLpIrr}{10.87}
\newcommand{\yaPref}{43.10}
\newcommand{\yaRest}{26.12}
\newcommand{\yaRocSix}{40.0}
\newcommand{\yaRocThree}{60.0}
\newcommand{\ybCarryEarly}{10}
\newcommand{\ybClaw}{6}
\newcommand{\ybFinal}{4}
\newcommand{\ybPv}{1.818}
\newcommand{\ycDur}{4.21}
\newcommand{\ydAlpha}{5.19}
\newcommand{\ydIrr}{12.47}
\newcommand{\ydPme}{1.231}
\newcommand{\yePme}{1.011}
\newcommand{\yePmeLow}{0.989}
\newcommand{\yfK}{1.292}
\newcommand{\yfMoic}{2.56}
\newcommand{\yfMoicLow}{1.06}
\newcommand{\ygD}{80}
\newcommand{\ygG}{142.5}
\newcommand{\ygM}{47.5}
\newcommand{\ygTot}{270}
\newcommand{\yhV}{160.0}
\newcommand{\yhVflat}{125.0}
\newcommand{\yhVhigh}{222.2}
\newcommand{\yiHp}{8.400}
\newcommand{\yiMe}{8.383}
\newcommand{\yjHam}{0.945}
\newcommand{\yjHp}{0.867}
\newcommand{\ykE}{8.25}
\newcommand{\ykV}{57.75}
\newcommand{\ylComb}{42.3}
\newcommand{\ylDloc}{23.1}
\newcommand{\ylStake}{5.19}
\newcommand{\ymLs}{53.7}
\newcommand{\ymPut}{17.8}
\newcommand{\ynFlat}{11.11}
\newcommand{\ynHigh}{8.57}
\newcommand{\ynLow}{10.20}

\raggedbottom % this chapter only: short pages stay short; \flushbottom is restored at the end
\begin{chapterrefs}
\setcounter{chapter}{10}
\chapter{Private Equity and Business Valuation}
\label{ch:pe}

A buyout fund raises money from pension plans, endowments and insurers, buys companies with that money
and with debt the companies themselves take on, runs them for a few years and sells them. Three questions follow,
and this chapter takes them in turn. How is the cash that comes back split between the investors and
the managers? How should an investor read the numbers a fund reports while it is still alive, and how
should it compare them with what the stock market would have paid? And what is a private company worth,
on the day the fund buys it, on the day it sells it, and on every reporting date between?

Chapters 1 and 2 have already done part of the work. Chapter~1 defined the IRR, the modified IRR, the
multiples DPI and TVPI, the Kaplan and Schoar public market equivalent and the J-curve\index{J-curve} of a single
fund. Chapter~2 took the fee contract apart: the management fee, the hurdle balance, the size of the
catch-up, carry as an option on a single exit, deal by deal carry and the clawback, and the
Takahashi and Alexander model of a fund's cash flows. This chapter builds on those results without
repeating them. It follows the waterfall of a whole fund through time, measures what deal by deal
carry is worth to the managers, puts a population of simulated funds and the \cvFunds{} partnerships in the
private equity portfolio of the California Public Employees' Retirement System (CalPERS) through the
performance measures, works out when debt raises a buyout's return, and then values a company: free
cash flow, the cost of capital, the terminal value, the market multiple, and the discounts that apply
to a stake that neither controls the company nor can be sold.

\begin{keybox}[How to read this chapter]
The numbers in the text are computed by six scripts in \texttt{code/ch11/} (\texttt{waterfall},
\texttt{performance}, \texttt{lbo}, \texttt{valuation}, \texttt{discounts} and \texttt{exercises}), and
the tests in \texttt{code/ch11/test\_*.py} check the library against simulation, finite differences and
identities that must hold exactly. An independent R script recomputes the waterfall tier by tier, the
IRRs with \texttt{uniroot}, the PME and direct alpha of the listing in \cref{sec11:perf}, the CalPERS
multiples and implied durations from the raw columns, the buyout model with its own loop, the cost of
capital and the DCF, and the two discounts by numerical integration. GNU Octave recomputes the
leverage proposition, the value driver\index{value driver}, the Miles and Ezzell cost of capital and the adjusted present
value by its own recursion, the lookback value by quadrature of the law of the maximum, and the
exercise answers.

The CalPERS fund list is the Private Equity Program Fund Performance Review as of 31 March 2026
\parencite{calpers2026}. Industry multiples, betas, costs of capital and returns on capital for 96 US
industries, and the implied equity risk premium from 1961 to 2025 (the file starts in 1960 with a Treasury
yield and no premium), are from the data sets of
Aswath Damodaran \parencite{damodaran2026}. The ten year Treasury yield is the September 2026 monthly
average from FRED \parencite{fredch11}. Neither the CalPERS nor the Damodaran files are redistributed
with the code; \texttt{scripts/pull\_ch11.py} downloads and rebuilds them. The fund, the deals, the
company and the exercise inputs are illustrations and are labelled as such where they appear.

Random draws come from NumPy's default generator with base seed \texttt{5279} and a stream number fixed
in each script. The seed is arbitrary; fixing it makes every run repeatable.

Notation. $C_t$ and $D_t$ are a fund's contributions and distributions at time $t$, $h$ the hurdle
rate, $c$ the carry and $k$ the catch-up rate. $I_t$ is the level of a public index, and $Z$ is a standard Brownian motion wherever it appears. In
\cref{sec11:lbo}, $\mathit{EV}$ is enterprise value, $\lambda$ debt as a share of what the
buyer pays, $i$ the interest rate, $\tau$ the tax rate and $\kappa$ the after-tax cost of debt compounded
over the holding period, and implied duration is written $\mathcal D$. In \cref{sec11:value}, $D$
and $E$ are the values of debt and equity, $\mathit{NOPAT}$ is
operating profit after tax, $\mathit{ROC}$ the return on new capital, $g$ a growth rate, $W$ the
weighted average cost of capital, $r_U$ the cost of capital of the unlevered firm, $r_D$ the cost of
debt and $L$ the ratio of debt to value. The letter $E$ has four uses: $\E$ is an expectation,
$E(t)$ in \cref{prop11:timing} is whole fund carry, $E$ in \cref{lem11:bridge} is EBITDA, and $E$ in
\cref{sec11:value} is the value of equity. Rates are annual and compounded once a year unless they are
marked as continuous. $\Phi$ and $\varphi$ are the standard normal distribution and density functions.
\end{keybox}

\FloatBarrier
% =====================================================================================
\section{How does a fund pay its managers?}
\label{sec11:fees}

A limited partnership commits capital at the start, and the general partner (GP) calls it as deals
are found, usually over an investment period of about five years. The limited partners (LPs) also fund a
management fee through capital calls, commonly a percentage of commitments during the investment
period and afterwards a percentage of the cost of the deals still held; the precise base varies by
fund. When a company
is sold, the proceeds pass through a \emph{distribution waterfall}\index{distribution waterfall}. Chapter~2 worked out each tier
for a fund with a single contribution and a single exit. A real fund calls and returns money many
times, and the tiers are tested again on each distribution date against everything that has happened
so far. That cumulative test is what the waterfall below formalises.

\begin{definition}[Whole fund waterfall]\label{def11:waterfall}
Let contributions $C_t\ge0$ and gross distributable proceeds $X_t\ge0$ occur at dates $t_0<t_1<\dots$.
The \emph{hurdle balance}\index{hurdle balance} accrues at the hurdle rate and rises with each contribution,
$B_{t_j}=B_{t_{j-1}}(1+h)^{t_j-t_{j-1}}+C_{t_j}$, and falls by whatever is paid to the LPs in the first two
tiers. On each date the proceeds $X_{t_j}$ are paid in order:
\begin{enumerate}[label=(\roman*)]
  \item to the LPs until all contributed capital has been returned;
  \item to the LPs until the hurdle balance, already reduced by the capital returned in (i), is zero,
        the \emph{preferred return}\index{preferred return};
  \item a \emph{catch-up}, in which a share $k>c$ of each dollar goes to the GP until the GP holds a share
        $c$ of the cumulative profit, $G=c\,(\text{LP distributions}+G-\text{contributions})$, with
        profit measured after fees and expenses, the convention of this chapter; agreements differ in
        how they treat fees, earlier carry and unrealised gains;
  \item the rest split $1-c$ to the LPs and $c$ to the GP.
\end{enumerate}
\end{definition}

The catch-up of tier (iii) has a closed size at each date. If the GP holds $G$ and the LPs have made
a cumulative profit $P$ when the tier opens, a payment of $y$ in the tier raises the GP's holding to
$G+ky$ and the profit to $P+(1-k)y$, and the tier closes when $G+ky=c\,(P+G+y)$, that is at
\begin{equation}\label{eq11:catchup}
  y^*=\frac{c\,(P+G)-G}{k-c}.
\end{equation}
This is the amount needed to reach the target; the tier pays $\min\bigl((y^*)^+,X\bigr)$ out of the
proceeds $X$ left for it. With a single exit, $G=0$ and \cref{eq11:catchup} is the catch-up size of
Chapter~2.

The listing below runs \cref{def11:waterfall} as a loop over distribution dates on the cash flows of
a fund that calls 100 at year 0 and distributes 50 at year 3 and 140 at year 6. At each
date the hurdle balance accrues, the tiers are tested again on the cumulative amounts, and the catch-up
takes the size that \cref{eq11:catchup} gives. It prints the payments by tier and the totals,
\wlGp{} to the GP and \wlLp{} to the LPs, so that the GP ends with a fifth of the profit of \wlProfit, as
\cref{prop11:share} below says.
\Needspace*{20\baselineskip}
\lstinputlisting[language=Python]{code/ch11/snippets/waterfall.py}

\begin{proposition}[What a full catch-up delivers]\label{prop11:share}
Suppose that, in the waterfall of \cref{def11:waterfall}, the catch-up tier has been completed on
some date and that no contribution is made after that date. Then from that date on the GP's
cumulative carry is exactly $c$ times the cumulative profit of the fund, distributions to both parties
less all contributions.
\end{proposition}

\begin{proof}
Assume, as the proposition does, that no later contribution, adjustment or clawback reopens tiers (i)
and (ii). When the catch-up tier closes, $G=c\,(P+G)$, where $P+G$ is the profit so far. On any later date
tiers (i) and (ii) pay nothing, because capital and hurdle were already cleared and no new
contribution has raised the hurdle balance, and the catch-up target is already met, so
\cref{eq11:catchup} gives $y^*=0$. Each later dollar is split by tier (iv) and raises $G$ by $c$ and
$P+G$ by one, so the equality $G=c\,(P+G)$ is preserved.
\end{proof}

The condition about contributions matters in practice. Management fees are called after the first
exits, and each such call lowers the fund's profit after the GP has already been paid a share of it.
The waterfall in the code tracks this and reports the excess the GP would owe back at the end; in the
fund below it is \wfClaw{}. Fees are called at the start of each year, the last of them in year eight on
the same date as the first carry payment and ahead of it in the waterfall, and none is called after it.

\paragraph{A stylised fund.}
Take a fund with commitments of \wfCommit{} that buys \wfDeals{} companies for \wfCost{} each over
its first five years and sells them between years three and nine at gross multiples of cost from 0 to
3.2. It charges \wfFeeInv\% of commitments a year for its first \wfInvYears{} years and \wfFeePost\% of
the cost of the deals still held after that, has a hurdle of \wfH\% and a carry of \wfCarry\% with a full
catch-up. All of this is an illustration. The LPs pay in \wfPaidIn, of which \wfFees{} is fees, and the
deals return \wfGross, a gross multiple of \wfGrossMoic{} on the capital invested. \Cref{fig11:waterfall}
shows where each year's distributions go. The exits of years three and five only return capital. The
preferred return is paid from year six to year eight, the catch-up of \wfCatch{} falls in year eight,
and the 80/20 split begins in the same year. Over the fund's life the LPs receive \wfLpTotal{} and the GP
\wfGpTotal, which is \wfGpShare\% of the profit of \wfProfit{} as \cref{prop11:share} requires. The
deals earn a gross IRR of \wfGrossIrr\%; the LPs earn \wfNetIrrFee\% after fees alone and \wfNetIrr\%
after fees and carry, with a TVPI of \wfNetTvpi.

\begin{figure}[!tb]
  \centering
  \includegraphics[width=0.94\textwidth]{ch11/fig01_waterfall}
  \caption{The stylised fund of \cref{sec11:fees}: contributions by year, including fees, below the
  axis, and each year's distributions split by the tier of \cref{def11:waterfall} that pays them.}
  \label{fig11:waterfall}
\end{figure}

\paragraph{From gross to net.}
Scaling every deal's multiple up and down traces the map from the gross IRR of the deals to the LPs'
net IRR (\cref{fig11:grossnet}). Fees cost the LPs a few points at every level of performance. For this
fund's cash flows and a full catch-up, carry costs nothing below the hurdle, then opens a nearly flat
stretch where the LPs' IRR stays at \wfH\% while the
gross IRR rises, which is the catch-up taking every marginal dollar, and then grows in proportion. At
a gross multiple of \wfGmLow{} the spread between gross and net IRR is \wfSpreadLow{} points, at
\wfGmMid{} it is \wfSpreadMid{} and at \wfGmHigh{} it is \wfSpreadHigh. The LPs clear the hurdle once
the deals return \wfMoicAtHurdle{} times their cost, a gross IRR of \wfGrossAtHurdle\%. The spread
between gross and net returns is therefore not a fixed cost: it rises with performance, and a spread
quoted for one fund says little about another, and other contracts give the map other shapes.

\begin{figure}[!tb]
  \centering
  \includegraphics[width=0.84\textwidth]{ch11/fig02_gross_net}
  \caption{IRR to the LPs against the gross IRR of the deals, when every deal multiple of the stylised
  fund is scaled by a common factor. The flat stretch at the hurdle is the catch-up.}
  \label{fig11:grossnet}
\end{figure}

\Cref{fig11:catchup} redraws the stretch around the hurdle for catch-up rates $k$ of 1, 0.8 and 0.5 and
with no catch-up. A slower catch-up shares the proceeds above the hurdle with the LPs, so the flat
stretch tilts upward and lengthens. The GP holds its full \wfCarry\% of the profit once the deals return
\ckFullOne{} times their cost with a full catch-up, \ckFullEight{} times at 0.8 and \ckFullHalf{} times at
0.5, and above those multiples the three contracts pay the same. With no catch-up the GP never reaches
its share: for this fund it ends with \ckShareZero\% of the profit, and at a gross multiple of
\wfGmMid{} the LPs earn \ckZeroMid\% instead of \wfNiMid\%.

\begin{figure}[!tb]
  \centering
  \includegraphics[width=0.84\textwidth]{ch11/fig15_catchup}
  \caption{IRR to the LPs against the gross IRR of the deals for the stylised fund of
  \cref{fig11:grossnet}, around the hurdle of \wfH\%, with catch-up rates $k$ of 1 (the contract of the
  text), 0.8 and 0.5 and with no catch-up.}
  \label{fig11:catchup}
\end{figure}

\subsection{What deal by deal carry is worth to the managers}

Under deal by deal (often called American) terms, carry is paid on each profitable exit as it happens,
without waiting for the fund to return all the capital it has called. A clawback at the end reconciles
the total with what whole fund (European) terms would have paid. Chapter~2 compared the two totals
when the clawback cannot be collected in full. Here the clawback is assumed to be paid in full, and
the question is what the earlier payments are worth.

\begin{proposition}[Deal by deal carry is an interest-free loan]\label{prop11:timing}
Let deal $i$ cost $K_i$, paid at $u_i$, and return $P_i$ at $s_i>u_i$, with no hurdle and carry $c$.
Write $E(t)=c\,\bigl(\sum_{s_i\le t}P_i-\sum_{u_i\le t}K_i\bigr)^+$ for a benchmark path of cumulative whole
fund carry by date $t$, an entitlement that need not match what a real whole fund agreement pays in the
interim, and $A(t)=c\sum_{s_i\le t}(P_i-K_i)^+$ for the carry paid deal by deal, until the last exit $T$,
where a full clawback sets $A(T)=E(T)$. Then $A(t)\ge E(t)$ for every $t$, and at any continuously compounded discount
rate $r\ge0$ the present value of the deal by deal payments exceeds that of the whole fund payments by
\begin{equation}\label{eq11:timing}
  \sum_j\bigl(A(t_j)-E(t_j)\bigr)\bigl(e^{-rt_j}-e^{-rt_{j+1}}\bigr)\;\ge\;0,
\end{equation}
where $t_1<t_2<\dots<t_n$ are the dates on which either path changes and $t_{n+1}=T$.
\end{proposition}

\begin{proof}
For $t<T$, the realised deals contribute $\sum_{s_i\le t}(P_i-K_i)$ to the bracket in $E(t)$, and every
deal called but not yet sold subtracts its cost, so the bracket is at most
$\sum_{s_i\le t}(P_i-K_i)\le\sum_{s_i\le t}(P_i-K_i)^+$. Both sides are then multiplied by $c\ge0$, and
$x\mapsto x^+$ is increasing, so $E(t)\le A(t)$. For the present values, write $\Delta_j$ for the
change in $A-E$ at $t_j$. The difference in present value is $\sum_j\Delta_je^{-rt_j}$, and summing by
parts with $A-E$ equal to zero before $t_1$ and at $T$ gives \cref{eq11:timing}. Each term is a
nonnegative difference $A-E$ times a nonnegative difference of discount factors. The argument compares
two cumulative paths that end at the same amount, clawback included; it does not show that every
interim value of $E$ is carry actually paid.
\end{proof}

\Cref{fig11:carrypaths} draws the two paths for a toy fund of three deals of 10 each, all bought at year 0, with no hurdle.
The first deal is sold for \ctPOne{} in year \ctSOne, the second for \ctPTwo{} in year \ctSTwo, and the
third, a loser, for \ctPThree{} in year \ctSThree. Deal by deal, the GP takes a fifth of each gain as it is
realised, \ctAPeak{} by year \ctSTwo; the whole fund rule pays only once the realised proceeds exceed the
capital called. The loser is sold last, so the GP returns \ctClaw{} at the clawback and both rules end at
\ctFinal. The shaded area is the carry held early, and discounted at a continuously compounded \ctR\% it is worth \ctPvGap{} to the
GP, the sum in \cref{eq11:timing}.

\begin{figure}[!tb]
  \centering
  \includegraphics[width=0.8\textwidth]{ch11/fig16_carrypaths}
  \caption{Cumulative carry of a toy fund of three deals, the loser sold last, under deal by deal carry
  with a full clawback, $A(t)$, and under whole fund carry, $E(t)$, with no hurdle. The shaded gap is
  carry paid early.}
  \label{fig11:carrypaths}
\end{figure}

The proof shows what the GP gains: on every interval it holds carry it may later have to return, and
it holds it free of interest. A simulation puts a size on the gain. Each of \wfTmFunds{} funds buys
\wfTmDeals{} deals of equal cost at dates spread evenly over its first four years, holds each for a
random time with a mean of four years, and sells it at a lognormal multiple with drift \wfTmMu\% and
volatility \wfTmSig\% a year. The deal by deal path is worth \wfTmExMed\% of the capital invested more
than the whole fund path at a continuously compounded discount rate of \wfTmR\% in the median fund, and \wfTmExNine\% in the
ninetieth percentile (\cref{fig11:timing}); across all funds the average gain is \wfTmRatio\% of the average present
value of the whole fund carry. The clawback is no formality: \wfTmClawFreq\% of the funds end with
the GP owing some carry back, and in the median such fund the amount is \wfTmClawMed\% of the carry
it finally keeps. Whether that money comes back depends on the clawback's terms, the taxes the GP has
already paid on it and the GP's solvency, which is the subject of the clawback section in Chapter~2.

\begin{figure}[!tb]
  \centering
  \includegraphics[width=0.84\textwidth]{ch11/fig03_timing}
  \caption{Present value at a continuously compounded \wfTmR\% of the deal by deal carry less the whole fund carry, in percent of
  the capital invested, across \wfTmFunds{} simulated funds with a full clawback and no hurdle.}
  \label{fig11:timing}
\end{figure}

\FloatBarrier
% =====================================================================================
\section{What do a fund's numbers measure?}
\label{sec11:perf}

An LP sees a fund through its quarterly report: cash paid in, cash paid out, and the GP's estimate of
the value of what is left, the net asset value (NAV). Chapter~1 turned these into an interim IRR, the
multiples DPI and TVPI, and the Kaplan and Schoar PME. This section asks how those numbers behave
across many funds, what an IRR and a multiple reported together imply about timing, and what the PME
does and does not adjust for.

\subsection{A population of simulated funds}

Each simulated fund lives \jcLife{} years and buys \jcDealsPer{} companies at quarters drawn evenly
over its first five years, investing 90\% of its commitments; the management fees of the first five
years use the other 10\%, and the fees after that are called on top of the commitments. A deal is held for at least a year and
then exits with probability 1/14 a quarter, so its mean holding period would be four and a half years
without the end of the fund; any deal still held at the end is sold then, which shortens the actual
holdings. The fund has its own path of a public index $I_t$, a geometric Brownian motion with drift
$\mu_I=\jcMuM\%$, its expected return, and volatility $\sigma_I=\jcSigM\%$ a year. A deal bought at $u$ for
$K$ is worth at $t$
\begin{equation}\label{eq11:deal}
  V_t=K\exp\Bigl(\bigl(\alpha-\tfrac12\sigma_e^2\bigr)(t-u)+\beta\ln\frac{I_t}{I_u}
  -\tfrac12\beta(\beta-1)\sigma_I^2(t-u)+\sigma_e\,(Z_t-Z_u)\Bigr),
\end{equation}
with $Z$ a Brownian motion of its own, $\alpha=\jcAlpha\%$, $\beta=\jcBeta$ and $\sigma_e=\jcSigE\%$. The
correction term makes the expected growth rate of $V$ exactly $\alpha+\beta\mu_I$, so $\alpha$ is the
return above what an unfinanced position of $\beta$ in the index would earn; the simulation has no
risk-free rate, so the cost of financing such a position does not enter. The deal's debt is not
modelled on its own but folded into $\beta$ and $\sigma_e$, although in real buyouts the amount of debt
follows conditions in the credit markets \parencite{axelson2013}. Fees are those of the stylised fund.
The GP marks each deal at cost for its first year and at $V_t$ afterwards, and the reported NAV is
net of a carry of 20\% of any profit to date, distributions plus NAV less contributions. There are
\jcFunds{} funds.

\Cref{fig11:jcurve} shows the J-curve of Chapter~1 as a distribution. After one year the median
interim IRR is \jcIrrOne\%, driven mainly by fees while most deals are still marked at cost in their first year. After two years
it is \jcIrrTwo\%, with a tenth of funds below \jcIrrLoTwo\% and a tenth above \jcIrrHiTwo\%, and the
median turns positive after \jcCross{} years. The median fund ends with an IRR of \jcIrrEnd\% and a
TVPI of \jcTvpiEnd. DPI lags TVPI by years: after four years the median fund has returned \jcDpiFour{}
of its paid in capital in cash while reporting a TVPI of \jcTvpiFour.

\begin{figure}[!tb]
  \centering
  \includegraphics[width=0.94\textwidth]{ch11/fig04_jcurve}
  \caption{(a) Interim net IRR of \jcFunds{} simulated funds by age: median and the band from the
  tenth to the ninetieth percentile. (b) Median TVPI and DPI by age.}
  \label{fig11:jcurve}
\end{figure}

How much does an early IRR say about the final one? The rank correlation between the interim IRR
and the final IRR is \jcRankThree{} after three years, \jcRankFive{} after five and \jcRankEight{} after
eight. Of the funds that end in the top quartile, \jcTopThree\% were in the top quartile at three
years, \jcTopFive\% at five and \jcTopEight\% at eight. In this model the record of a fund three years old, about the
age at which a successor fund is often marketed, predicts the final ranking only loosely: fewer than
half of the eventual top quartile is in the top quartile at that age.

\Cref{fig11:transition} gives the whole picture at three and five years. Each row is a quartile of the
final IRR and shows where those funds stood at the earlier age; the top right cells are the \jcTopThree\%
and \jcTopFive\% above. Of the funds that end in the top quartile, \trTopBottomThree\% were in the bottom
quartile at three years and \trTopBottomFive\% at five, and of those that end in the bottom quartile,
\trBottomTopThree\% and \trBottomTopFive\% were in the top quartile. At three years a fund can still end
almost anywhere; by five, large moves are rare.

\begin{figure}[!tb]
  \centering
  \includegraphics[width=0.92\textwidth]{ch11/fig17_transition}
  \caption{Simulated funds: the share of the funds in each quartile of the final IRR that were in each
  quartile of the interim IRR at three and five years, in percent, rounded so that each row adds to 100. The
  funds are those of \cref{fig11:jcurve}.}
  \label{fig11:transition}
\end{figure}

\subsection{An IRR and a multiple together}

A fund reports an IRR and a TVPI side by side, and the pair carries information that neither number
has alone. If a single payment of 1 grew into a single payment of $M$ after $\mathcal D$ years at the rate $r$,
then $M=(1+r)^{\mathcal D}$. For a real fund, define the \emph{implied duration}\index{implied duration} by the same relation,
$\mathcal D=\ln\mathrm{TVPI}/\ln(1+r)$, with the NAV counted as a distribution on the reporting date.

\begin{proposition}[Implied duration]\label{prop11:duration}
Let a fund make contributions $C_i>0$ at times $t_i$ and distributions $D_j>0$ at times $s_j$, the last
of which may be the NAV, and let its IRR $r$ be positive. Then
\begin{equation}\label{eq11:duration}
  \min_js_j-\max_it_i\;\le\;\frac{\ln\mathrm{TVPI}}{\ln(1+r)}\;\le\;\max_js_j-\min_it_i .
\end{equation}
\end{proposition}

\begin{proof}
Write $\rho=\ln(1+r)>0$. The IRR equates the present values, $\sum_jD_je^{-\rho s_j}=\sum_iC_ie^{-\rho t_i}$.
Each discount factor on the left is at least $e^{-\rho\max s_j}$ and each on the right at most
$e^{-\rho\min t_i}$, so $e^{-\rho\max s_j}\sum_jD_j\le e^{-\rho\min t_i}\sum_iC_i$, that is
$\mathrm{TVPI}\le e^{\rho(\max s_j-\min t_i)}$. Taking logarithms and dividing by $\rho$ gives the upper
bound. Bounding the left side above by $e^{-\rho\min s_j}\sum_jD_j$ and the right side below by
$e^{-\rho\max t_i}\sum_iC_i$ gives the lower bound in the same way.
\end{proof}

The upper bound is the useful one. A fund with an IRR of 25\% and a TVPI of 1.2 has an implied
duration of $\ln1.2/\ln1.25=0.82$ years, a single period equivalent that is not the measured
average holding period of the fund's dollars; the high IRR is earned over an equivalent period of less
than a year. The bound is also a check on
data. A fund of vintage $v$ that reports as of the start of the second quarter of 2026 cannot have an
implied duration above $2026.25-v$ unless it called capital before 1 January of its vintage year or
the IRR and the multiple are inconsistent.

\Cref{fig11:duration} draws lines of equal implied duration in the plane of IRR and TVPI, from
$\mathrm{TVPI}=(1+r)^{\mathcal D}$. The example sits below the line of one year, at \duEx. The same reading
works the other way: a TVPI of 2 earned at an IRR of 10\% implies a duration of \duTen{} years, so at a
given multiple a lower IRR means the money was out for longer.

\begin{figure}[!tb]
  \centering
  \includegraphics[width=0.8\textwidth]{ch11/fig18_duration}
  \caption{Lines of equal implied duration $\mathcal D=\ln\mathrm{TVPI}/\ln(1+r)$, with the example of the
  text marked. The lines are analytic and carry no fund data.}
  \label{fig11:duration}
\end{figure}

\paragraph{The CalPERS portfolio.}
The CalPERS review lists \cvFunds{} active partnerships, from vintage \cvFirstVintage{} onwards, with
commitments of \cvCommitted{} billion dollars and a remaining value of \cvNav{} billion. For each it gives
cash contributed and distributed, total value and the net IRR, and it
flags the IRR as not meaningful for the \cvNm{} funds of vintage \cvNmFrom{} and later. The multiple it prints is total value over
cash contributed, which the code checks against the raw columns. Pooled by vintage, the TVPI is
\cvTvpiSeven{} for 2007, \cvTvpiFifteen{} for 2015, \cvTvpiNineteen{} for 2019 and \cvTvpiTwentyThree{}
for 2023, and the DPI \cvDpiSeven, \cvDpiFifteen, \cvDpiNineteen{} and \cvDpiTwentyThree{}
(\cref{fig11:calpers}a). The 2019 vintage has returned half its paid in capital in cash and the 2023
vintage almost none; their TVPIs are mostly the GPs' marks. The list holds active partnerships only,
so a fund of an early vintage that has been wound up is missing, and the \cvOldN{} funds of vintage
2012 or earlier that remain have a pooled TVPI of \cvOldTvpi{} with \cvOldRvpi\% of their total value
still unrealised.

\begin{figure}[!tb]
  \centering
  \includegraphics[width=0.94\textwidth]{ch11/fig06_calpers}
  \caption{The CalPERS private equity partnerships as of 31 March 2026. (a) TVPI and DPI pooled by
  vintage, for vintages with at least five funds. (b) Implied duration of each fund with a meaningful
  IRR above zero, against the bound of \cref{prop11:duration} measured from 1 January of the vintage.}
  \label{fig11:calpers}
\end{figure}

For the \cvDurN{} funds with a meaningful positive IRR, the median implied duration is \cvDurMed{}
years (\cref{fig11:calpers}b). No fund lies above the bound measured from 1 January of its vintage, and
the largest implied duration is \cvDurRatioMax\% of the time since its vintage began, so the reported
IRRs and multiples are consistent in the sense of \cref{prop11:duration}. The median fund's implied duration is \cvDurShareMax\% of the time since its
vintage began, which is consistent with contributions and distributions spread over a shorter
effective interval than the fund's age; the implied duration alone does not identify their timing. The funds of 2005 to 2008 have a median implied duration of \cvDurPre{} years, those
of 2015 to 2018 \cvDurMid.

\subsection{Comparing with the index}

The PME of \textcite{kaplan2005}, defined in Chapter~1, carries every cash flow to the end date with
the index and divides compounded distributions by compounded contributions. \textcite{gredil2023}
proposed the rate that goes with it.

\begin{definition}[Direct alpha]\label{def11:alpha}
Carry each contribution and distribution of a fund to the end date $T$ with the index,
$\tilde C_t=C_tI_T/I_t$ and $\tilde D_t=D_tI_T/I_t$, and count the NAV at $T$ once, as a final distribution
that is not also part of $\tilde D_T$. With time measured from the first cash flow, $t=0$, the
\emph{direct alpha}\index{direct alpha} $a$ is the continuously compounded IRR of the carried flows, the root of
$\sum_t(\tilde D_t-\tilde C_t)e^{-at}+\mathrm{NAV}_Te^{-aT}=0$.
\end{definition}

\begin{proposition}[PME and direct alpha agree on the sign]\label{prop11:pme}
Suppose that there is a date $t^*$ such that every carried net flow $\tilde D_t-\tilde C_t$ before
$t^*$ is negative or zero and every one after $t^*$, including the NAV, is positive or zero, with at
least one nonzero flow strictly before $t^*$ and one strictly after it. Then the direct alpha exists and is unique, and $\mathrm{KS\text{-}PME}>1$ if and
only if $a>0$.
\end{proposition}

\begin{proof}
Let $f(a)=\sum_t(\tilde D_t-\tilde C_t)e^{-at}$, with the NAV included, and $\psi(a)=e^{at^*}f(a)$, which has
the same roots as $f$. In $\psi$, a flow at $t<t^*$ is multiplied by $e^{a(t^*-t)}$, which increases with
$a$, and the flow is negative; a flow at $t>t^*$ is multiplied by $e^{-a(t-t^*)}$, which decreases
with $a$, and the flow is positive. So $\psi$ is strictly decreasing, it tends to $-\infty$ as
$a\to\infty$ and to $+\infty$ as $a\to-\infty$, and it has exactly one root. Since $\psi$ decreases,
the root is positive if and only if $\psi(0)>0$. And $\psi(0)=f(0)=\sum\tilde D-\sum\tilde C$, which is
positive exactly when the KS-PME\index{Kaplan-Schoar}, $\sum\tilde D/\sum\tilde C$, exceeds one.
\end{proof}

The condition is sufficient and not necessary. In the simulated funds, fees are called after the
first distributions, so the condition fails for many of them, yet the PME and the direct alpha agree
on the sign for every one of the \jcFunds{} funds, and their rank correlation is \pkCorr. The listing
below computes both for a fund of five cash flows, short enough to retype.
\Needspace*{16\baselineskip}
\lstinputlisting[language=Python]{code/ch11/snippets/pme_alpha.py}
The fund has a TVPI of \pkSnipTvpi{} and an IRR of \pkSnipIrr\%, and the listing prints a KS-PME of
\pkSnipPme{} and a direct alpha of \pkSnipAlpha\% a year. Across the simulated funds the median PME is
\pkMed, only \pkAbove\% of funds beat the index, and the median direct alpha is \pkDaMed\% a year, even
though every deal's value is given an alpha of \jcAlpha\% a year in the model, before fees and carry: fees and carry take more than the alpha, and the
dispersion of outcomes leaves most funds below the index. The comparison is not beta-neutral either: the
simulation has no risk-free rate, so by \cref{prop11:beta} with $r=0$ a beta of \jcBeta{} adds
$(\beta-1)(\mu_I-\sigma_I^2)=\pkBetaBoost\%$ a year to each deal's expected PME exponent. Run on the same
draws with a beta of one, the funds have a median PME of \pkMedBetaOne{} and \pkAboveBetaOne\% of them beat the index.

\subsection{What the PME assumes about risk}

The PME prices a fund's cash flows as if they had the risk of the index. \textcite{sorensen2015} showed
that it is the valuation an investor with logarithmic utility who holds the index as wealth would put
on the flows, with stochastic discount factor $I_t/I_T$ between $t$ and $T$. Under that model the index's
expected return over the risk-free rate equals its variance, $\mu-r=\sigma^2$, and every fund with no
alpha has an expected PME of one, whatever its beta. When the market's premium differs from $\sigma^2$, a
fund whose beta is not one shows a PME away from one with no skill, and the bias has a closed form in
the simplest case.

\begin{proposition}[Beta and the PME]\label{prop11:beta}
Let the index be a geometric Brownian motion with drift $\mu$ and volatility $\sigma$, so that
$\ln(I_T/I_0)$ is normal with mean $m=(\mu-\frac12\sigma^2)T$ and variance $v=\sigma^2T$, and let $r$ be a
constant, continuously compounded risk-free rate. A fund makes one contribution of 1 at time 0 and one distribution at $T$ of
\[
  e^{\alpha T}\,e^{(1-\beta)rT-\frac12\beta(\beta-1)v}\,(I_T/I_0)^\beta,
\]
which is $e^{\alpha T}$ times the value of a position of $\beta$ in the index, rebalanced continuously and
financed at $r$. Then
\begin{equation}\label{eq11:betapme}
  \E\,\mathrm{KS\text{-}PME}=\exp\bigl(\alpha T+(\beta-1)(\mu-r-\sigma^2)T\bigr).
\end{equation}
With $\alpha=0$ the expected PME exceeds one if and only if $(\beta-1)(\mu-r-\sigma^2)>0$.
\end{proposition}

\begin{proof}
A position of $\beta$ in the index financed at $r$ has $dV/V=\beta\,dI/I+(1-\beta)r\,dt$, so by Itô's formula
$\ln(V_T/V_0)=\beta\ln(I_T/I_0)+(1-\beta)rT-\frac12\beta(\beta-1)\sigma^2T$, the stated distribution. The PME
is the distribution divided by the index return,
$e^{\alpha T+(1-\beta)rT-\frac12\beta(\beta-1)v}(I_T/I_0)^{\beta-1}$, which is lognormal with log mean
$\alpha T+(1-\beta)rT-\frac12\beta(\beta-1)v+(\beta-1)m$ and log variance $(\beta-1)^2v$. The mean of a
lognormal variable is the exponential of its log mean plus half its log variance, and with
$m=(\mu-\frac12\sigma^2)T$ the terms in $\sigma^2$ collect to $-(\beta-1)\sigma^2T$, which gives
\cref{eq11:betapme}. The last statement reads off the sign of the exponent.
\end{proof}

With the simulation's index parameters, $\mu=\pkMu\%$ and $\sigma=\pkSig\%$, but leverage financed at the Treasury yield of \wkRf\%, used here as a continuously
compounded rate, and $T=\pkT$ years, the premium $\mu-r=\pkPrem\%$ is below $\sigma^2=\pkVar\%$, and a fund with no alpha has
an expected PME of \pkBetaHalf{} at a beta of 0.5, \pkBetaOneHalf{} at 1.5 and \pkBetaTwo{} at 2; a
simulation at a beta of 2 gives \pkBetaSimTwo{} with standard error \pkBetaSimTwoSE{} (\cref{fig11:beta}).
The sign is the reverse of the simulation's in \cref{sec11:perf}, where the deals are unfinanced,
$r=0$, and $\mu_I-\sigma_I^2=\jcMuLessVar\%$ is positive.
Here high beta lowers the PME: the log of the expected PME at a beta of 1.5, divided by five, is
\pkFakeAlpha\% a year, a model quantity and not an estimate of alpha. With a premium of \pkErp\%, the
implied premium used for the cost of capital below, the same fund at a beta of 2 has an expected PME of
\pkBetaTwoErp. Buyout equity is levered, and many published estimates of its beta exceed one, so a PME
above one can be partly a reward for leverage, but only when the market pays more per unit of risk
than the logarithmic utility benchmark assumes.
\textcite{korteweg2016} generalise the PME to a discount factor that is exponentially affine in the
log index return, with parameters chosen so that it prices public assets, which corrects the bias when a
fund's risk loads on the index in that way.

\begin{figure}[!tb]
  \centering
  \includegraphics[width=0.84\textwidth]{ch11/fig05_beta}
  \caption{Expected KS-PME of a fund with no alpha, one contribution and one distribution five years
  later, against its beta on the index, from \cref{eq11:betapme}, with leverage financed at the
  Treasury yield of \wkRf\%, and simulated means with two standard errors.}
  \label{fig11:beta}
\end{figure}

\FloatBarrier
% =====================================================================================
\section{When does debt raise a buyout's return?}
\label{sec11:lbo}

In a leveraged buyout\index{leveraged buyout} the fund pays for a company partly with equity and partly with debt raised
against the company's own cash flows, and the company repays the debt from those cash flows. Without
taxes and frictions, \textcite{modigliani1958} showed that the mix of debt and equity does not change
the value of the firm; with tax deductible interest the debt adds the value of its tax shields
\parencite{modigliani1963}. Leverage does change the return on the equity, in a way that has a short
exact form.

\begin{proposition}[Leverage and the equity multiple]\label{prop11:lbo}
A buyer pays $P_0$ for a company, fees included. Without debt, the company would at $T$ be worth
$\mathit{EV}_T$ plus its free cash flows $F_t$, each held to $T$ at the rate $j=i(1-\tau)$, and its
\emph{asset multiple}\index{asset multiple} is $U=\bigl(\mathit{EV}_T+\sum_tF_t(1+j)^{T-t}\bigr)/P_0$. Suppose instead that debt
$\lambda P_0$ is borrowed at the rate $i$, that interest is deductible at the tax rate $\tau$ against
taxable income that stays positive, and that all cash after interest either repays debt or earns $j$.
If the equity is worth something at $T$, its multiple is
\begin{equation}\label{eq11:lbo}
  M_E=\frac{U-\lambda \kappa }{1-\lambda},\qquad \kappa=(1+j)^T,\qquad\text{so}\qquad
  M_E-U=\frac{\lambda}{1-\lambda}\,(U-\kappa).
\end{equation}
Debt raises the equity multiple if and only if $U>\kappa$, and for a fixed $U>\kappa$ the gain grows without bound
as $\lambda$ rises to one; solvency and the limits a lender would set are outside the model.
\end{proposition}

\begin{proof}
Let $N_t$ denote debt less cash, which may turn negative. In this simplified model, borrowing and
cash balances carry the same pretax rate $i$, and interest expense and interest income receive symmetric
tax treatment at the rate $\tau$. Consequently net debt accrues at $j=i(1-\tau)$, before the free cash
flow $F_t$ the company would have generated without debt is applied. So $N_t=N_{t-1}(1+j)-F_t$ with $N_0=\lambda P_0$, and by induction
$N_T=\lambda P_0\kappa-\sum_tF_t(1+j)^{T-t}$. The equity is worth $\mathit{EV}_T-N_T=P_0U-\lambda P_0\kappa$ and cost
$(1-\lambda)P_0$, which gives \cref{eq11:lbo}; the difference $M_E-U$ follows by subtraction.
\end{proof}

The factor $\kappa$ is the debt's cost after tax, compounded over the holding period. At $i=\lxI\%$,
$\tau=\lxTau\%$ and $T=\lxT$ years it is \lxK, so leverage helps exactly when the company would have
grown its value at more than 6\% a year on its own. \Cref{fig11:leverage} draws the equity multiple
against $\lambda$. At $\lambda=\lxLamEx\%$, an asset multiple of \lxUHi{} becomes an equity multiple of
\lxMHi, and an asset multiple of \lxULo{} becomes \lxMLo: the same debt that turns a good deal into a
very good one turns a mediocre deal into a loss. The proposition holds deal by deal, so it also says
that leverage cannot be judged by its average effect: a fund's leverage decision is a bet that its
deals will clear $\kappa$.

\begin{figure}[!tb]
  \centering
  \includegraphics[width=0.84\textwidth]{ch11/fig07_leverage}
  \caption{Equity multiple against debt as a share of the price, from \cref{eq11:lbo}, for four asset
  multiples $U$, with $i=\lxI\%$, $\tau=\lxTau\%$ and $T=\lxT$ years.}
  \label{fig11:leverage}
\end{figure}

\subsection{A five year model}

\Cref{tab11:lbo} sets out a buyout of a stylised company with revenue of \lxRev{} million and EBITDA of
\lxEz{} million, bought at \lxMz{} times EBITDA, an enterprise value of \lxEvz{} million, with fees of
\lxFees\% of that value. The debt is \lxLev{} times EBITDA, \lxDz{} million at \lxRate\%, and every dollar of
cash after interest, tax, capital spending of \lxCapex\% of revenue and working capital of \lxNwc\% of the
growth in revenue repays it. Revenue grows by 6\%, 6\%, 5\%, 5\% and 4\%, the
EBITDA margin rises by half a point a year from 20\% to 22\% in the fourth year and stays there, and tax
is 25\% of EBITDA less depreciation of 3\% of revenue and less interest. The equity cheque is \lxEqz{} million. After five years EBITDA is
\lxEt{} million; sold at \lxMt{} times, the company is worth \lxEvt{} million, the debt has fallen to
\lxNdt{} million, \lxLevExit{} times EBITDA, and the equity is worth \lxEqt{} million, a multiple of
\lxMoic{} and an IRR of \lxIrr\%. The same deal without debt returns \lxMoicU{} times, an IRR of
\lxIrrU\%. The code checks that \cref{eq11:lbo}, with $U=\lxUb$ and $\lambda=\lxLamB\%$ of the price
including fees, gives the levered multiple of the model to rounding error.

\begin{table}[!tb]
  \centering
  \small
  \caption{The stylised buyout, in millions of dollars. Free cash flow is after interest, tax,
  capital spending and working capital, and all of it repays debt.}
  \label{tab11:lbo}
  \begin{tabular}{@{}lrrrrr@{}}
    \toprule
    Year & 1 & 2 & 3 & 4 & 5\\
    \midrule
    Revenue & \lxTRevOne & \lxTRevTwo & \lxTRevThree & \lxTRevFour & \lxTRevFive\\
    EBITDA & \lxTEbitdaOne & \lxTEbitdaTwo & \lxTEbitdaThree & \lxTEbitdaFour & \lxTEbitdaFive\\
    Interest & \lxTInterestOne & \lxTInterestTwo & \lxTInterestThree & \lxTInterestFour & \lxTInterestFive\\
    Tax & \lxTTaxOne & \lxTTaxTwo & \lxTTaxThree & \lxTTaxFour & \lxTTaxFive\\
    Capital spending & \lxTCapexOne & \lxTCapexTwo & \lxTCapexThree & \lxTCapexFour & \lxTCapexFive\\
    Working capital & \lxTDnwcOne & \lxTDnwcTwo & \lxTDnwcThree & \lxTDnwcFour & \lxTDnwcFive\\
    Free cash flow & \lxTFcfOne & \lxTFcfTwo & \lxTFcfThree & \lxTFcfFour & \lxTFcfFive\\
    Debt at year end & \lxTDebtOne & \lxTDebtTwo & \lxTDebtThree & \lxTDebtFour & \lxTDebtFive\\
    \bottomrule
  \end{tabular}
\end{table}

The listing below is the model of \cref{tab11:lbo} in about twenty lines. It prints the exit EBITDA,
value, net debt and equity, the multiple of \lxMoic{} and the IRR of \lxIrr\%, and the same deal without
debt, \lxMoicU{} times and \lxIrrU\%. Changing the debt multiple in the call, or the exit multiple inside
the function, is a quick way to check \cref{eq11:lbo} by hand.
\Needspace*{24\baselineskip}
\lstinputlisting[language=Python]{code/ch11/snippets/lbo.py}

\subsection{Where the gain came from}

Sponsors explain a deal's gain as EBITDA growth, a higher multiple at exit and debt repaid. The first
two interact, because value is their product, and a split that credits growth at the entry multiple
gives a different answer from one that credits it at the exit multiple. The symmetric split resolves
this.

\begin{lemma}[Splitting a product]\label{lem11:bridge}
For $V=mE$, with $\bar m=\frac12(m_0+m_T)$ and $\bar E=\frac12(E_0+E_T)$,
\[
  m_TE_T-m_0E_0=\bar m\,(E_T-E_0)+\bar E\,(m_T-m_0),
\]
and each term is the average of the two contributions that factor makes when it is changed first and
when it is changed second, which is its Shapley\index{Shapley} value in the two player game.
\end{lemma}

\begin{proof}
Expanding the right side gives $\frac12(m_0E_T-m_0E_0+m_TE_T-m_TE_0+E_0m_T-E_0m_0+E_Tm_T-E_Tm_0)$, in which
the cross terms cancel and $m_TE_T-m_0E_0$ remains. Changing $E$ first contributes $m_0(E_T-E_0)$ and
changing it second $m_T(E_T-E_0)$; their average is $\bar m(E_T-E_0)$, and likewise for $m$.
\end{proof}

The equity gain of the stylised deal, \lxGain{} million, splits into \lxBrG{} from EBITDA growth,
\lxBrM{} from the higher multiple and \lxBrD{} from debt repaid, less \lxBrF{} of fees
(\cref{fig11:bridge}). Crediting growth at the entry multiple would give it \lxGfirst{} and at the exit
multiple \lxMfirst. The debt term is \lxBrDs\% of the gain in the bridge, an accounting contribution and not a
separate source of value: the cash that repaid the debt came from the same operations that grew EBITDA, and an unlevered
owner would have held it as cash. What the bridge does show is how much the deal depended on the exit
multiple, which no operating plan controls; here it is \lxBrMs\%.

\begin{figure}[!tb]
  \centering
  \includegraphics[width=0.84\textwidth]{ch11/fig08_bridge}
  \caption{Bridge of the stylised deal's equity from entry to exit, in millions, with the interaction of
  EBITDA and the multiple split by \cref{lem11:bridge}.}
  \label{fig11:bridge}
\end{figure}

\Cref{fig11:grid} varies the entry and exit multiples between \lxGridMin{} and \lxGridMax, keeping the
debt at \lxLev{} times EBITDA. Buying at \lxGridMin{} and selling at \lxGridMin{} returns \lxGridLoLo\% a
year, and buying and selling at \lxGridMax{} only \lxGridHiHi\%: at the same multiple a cheaper entry
means that the fixed debt pays a larger share of the price. Buying at \lxGridMax{} and selling at
\lxGridMin{} still returns \lxGridHiLo\%, because EBITDA growth and repayment carry the deal, and the
reverse returns \lxGridLoHi\%. \textcite{axelson2013} found that buyout leverage follows the cost of debt
in the credit markets more closely than the characteristics of the companies bought, and that more
leverage goes with higher prices paid; in this grid, which holds the debt fixed, the high entry multiples of
the top rows are where returns at any given exit multiple are lowest.

\begin{figure}[!tb]
  \centering
  \includegraphics[width=0.78\textwidth]{ch11/fig09_grid}
  \caption{Equity IRR of the stylised buyout, in percent, by entry and exit EV/EBITDA, with debt of
  \lxLev{} times EBITDA at entry.}
  \label{fig11:grid}
\end{figure}

\paragraph{Uncertain outcomes.}
A simulation of \lxMcN{} deals draws revenue growth each year from a normal law with mean 5\% and
standard deviation 6\%, sets the margin
each year to 21\% plus 1.5 points times an independent normal draw plus a third of a point times the sum
of the standardised growth shocks so far, held between 5\% and 40\%, and draws the
exit multiple from a lognormal law with mean \lxMz{} and log standard deviation 18\%. With debt of
\lxLev{} times EBITDA the median IRR rises from \lxMcMedU\% to \lxMcMedL\% and its standard deviation
from \lxMcSdU{} to \lxMcSdL{} points (\cref{fig11:lbomc}). The share of deals that lose money rises from
\lxMcLossU\% to \lxMcLossL\%, and the fifth percentile of the IRR falls from \lxMcQU\% to \lxMcQL\%. Debt
raised the multiple in \lxMcBeat\% of the deals, which are exactly the \lxMcUAbove\% in which the asset
multiple exceeded $\kappa$, as \cref{prop11:lbo} says. The simulation has no default: a revolver absorbs any
shortfall, with no covenant, lender intervention or credit limit, so its downside figures understate
the risks that come with financial distress. \textcite{kaplan1989} documented large increases in operating
cash flow after the management buyouts of the 1980s, the operating side of the bridge, and
\textcite{harris2014} found that US buyout funds beat the S\&P 500 on a PME basis in most vintages they
studied. That result is consistent with strong buyout performance relative to public markets, although
\cref{prop11:beta} shows that a PME can also reflect systematic market exposure; the direction and
magnitude of any beta effect must be estimated for the actual funds and periods.

\begin{figure}[!tb]
  \centering
  \includegraphics[width=0.84\textwidth]{ch11/fig10_lbo_mc}
  \caption{Equity IRR over five years of \lxMcN{} simulated deals, without debt and with debt of
  \lxLev{} times EBITDA at entry.}
  \label{fig11:lbomc}
\end{figure}

\subsection{A cap on deductible interest}
\label{sec11:intcap}

The model so far deducts every dollar of interest, which is the assumption behind the factor
$1+i(1-\tau)$ in \cref{prop11:lbo}. Many tax systems now limit the deduction. The OECD recommended a
fixed ratio rule that caps net interest at between 10\% and 30\% of a tax measure of
EBITDA\index{interest deductibility cap} \parencite{oecd2015}. Canada adopted 30\% for taxation years
beginning on or after 1 January 2024 \parencite{ita182}, and interest denied in one year is carried
forward and can be deducted in a later year in which the company has room under the cap. The code
applies a cap of \lxCap\% of EBITDA, with EBITDA standing in for the tax measure, and carries the
excess forward.

In the base deal the first year's interest is \lxCapIntOne{} million, \lxCapIntShareOne\% of that
year's EBITDA, so only \lxCapDedOne{} million is deducted. The cap binds in the first \lxCapBinds{}
years, the amount carried forward peaks at \lxCapCarryPeak{} million, and the sweep pays the debt down
fast enough that the carryforward is used before exit. The cap therefore moves about \lxCapTaxEarly{}
million of tax into the early years without changing much the total paid, and the equity multiple
falls only from \lxMoicThree{} to \lxCapMoic, with the IRR lower by less than two basis points.
Without a carryforward the denied interest would be lost for good, the multiple would be
\lxCapNoCfMoic{} and the IRR \lxCapNoCfIrrDropBp{} basis points lower than without the cap. In this deal the cap
stops binding in the first year once debt at entry is below \lxCapLevFree{} times EBITDA.

The effect depends on how long the cap binds. A deal with seven times EBITDA of debt at 9\% exceeds the
cap in all \lxCapHvBinds{} years and still has \lxCapHvCarry{} million carried forward at exit, which
is worth something to a buyer only if the buyer can use it. Its IRR falls from \lxCapHvIrrNo\% to
\lxCapHvIrr\%, \lxCapHvIrrDropBp{} basis points. The cap also breaks the identity of
\cref{prop11:lbo}, since part of the interest is paid without a shield; the code checks that the
proposition overstates the capped multiple. The rules themselves have more parts than this sketch,
including a group ratio election and an exclusion for small entities, and the tax measure of EBITDA
differs from the accounting one.

\FloatBarrier
% =====================================================================================
\section{What is a private company worth?}
\label{sec11:value}

A buyout fund values a company when it bids, every quarter while it owns it, and when it sells. The
company has no market price, so the value is built from its cash flows, its cost of capital and the
prices paid for comparable companies. The standard approach is the discounted free cash flow to the
firm, and the cash flow it discounts is defined first.

\begin{definition}[Free cash flow to the firm]\label{def11:fcff}
The \emph{free cash flow to the firm} in year $t$ is $F_t=\mathit{NOPAT}_t-R_t$, where
$\mathit{NOPAT}_t=\mathit{EBIT}_t(1-\tau)$ is operating profit after tax at the marginal rate, as if the firm
had no debt, and the \emph{reinvestment} $R_t$ is capital spending less depreciation plus the increase in
working capital. The \emph{reinvestment rate}\index{reinvestment rate} is $R_t/\mathit{NOPAT}_t$, and the \emph{return on new capital}\index{return on new capital}
is the extra NOPAT that a dollar of reinvestment produces a year later.
\end{definition}

\subsection{The terminal value and the value driver}

A forecast runs for a few years and ends in a terminal value for everything after it. The growing
perpetuity of \textcite{gordon1956} is the usual choice, and it is usually applied to the last cash
flow of the forecast. That cash flow depends on how much the company reinvests, and the growth it
assumes has to be paid for by reinvestment. The value driver formula makes the link explicit.

\begin{proposition}[Value driver]\label{prop11:driver}
Suppose that from year $T+1$ on a firm has positive NOPAT, earns a stable return $\mathit{ROC}>0$ on new
capital, reinvests a constant share of NOPAT, and grows at $g<W$ forever. Then the reinvestment rate is $g/\mathit{ROC}$ and
\begin{equation}\label{eq11:driver}
  \mathit{TV}_T=\frac{\mathit{NOPAT}_{T+1}\,(1-g/\mathit{ROC})}{W-g},\qquad
  \frac{\partial}{\partial g}\,\frac{1-g/\mathit{ROC}}{W-g}=\frac{1-W/\mathit{ROC}}{(W-g)^2}.
\end{equation}
The derivative holds $\mathit{NOPAT}_{T+1}$ fixed. Growth adds value if and only if $\mathit{ROC}>W$; at $\mathit{ROC}=W$ the terminal value is
$\mathit{NOPAT}_{T+1}/W$ whatever the growth.
\end{proposition}

\begin{proof}
If a share $b$ of NOPAT is reinvested each year at the return $\mathit{ROC}$, next year's NOPAT exceeds
this year's by $b\,\mathit{ROC}\,\mathit{NOPAT}$, so $g=b\,\mathit{ROC}$ and $b=g/\mathit{ROC}$. The free cash flow is
$(1-b)\mathit{NOPAT}$ and grows at $g$, and the growing perpetuity gives the first formula. Differentiating
the quotient, the numerator of the derivative is $-(W-g)/\mathit{ROC}+(1-g/\mathit{ROC})=1-W/\mathit{ROC}$, which is
positive exactly when $\mathit{ROC}>W$. At $\mathit{ROC}=W$ the quotient is $(1-g/W)/(W-g)=1/W$.
\end{proof}

\Cref{fig11:driver} draws the multiple against $g$ at the manufacturer's cost of capital of \wkWme\%,
found below. Every curve starts at $1/W=\drZero$ when nothing is reinvested. At 5\% growth the multiple
is \drFiveTwenty{} with a return on new capital of 20\%, \drFiveTwelve{} at 12\% and only \drFiveSix{} at
6\%, below the multiple with no growth at all; with a return equal to $W$ the line is flat.

\begin{figure}[!tb]
  \centering
  \includegraphics[width=0.82\textwidth]{ch11/fig19_driver}
  \caption{Terminal value over next year's NOPAT from \cref{eq11:driver} against growth, at
  $W=\wkWme\%$, for four returns on new capital; the dot is the manufacturer's terminal value.}
  \label{fig11:driver}
\end{figure}

The proposition is a discipline on the inputs. A terminal growth rate chosen alone, without a return
on capital, leaves the reinvestment free. Letting $\mathit{ROC}$ grow without bound gives
$\mathit{NOPAT}_{T+1}/(W-g)$, growth that costs nothing; the usual perpetuity $F_{T+1}/(W-g)$ agrees with
the value driver only if $F_{T+1}$ already deducts the reinvestment $(g/\mathit{ROC})\,\mathit{NOPAT}_{T+1}$
that the growth requires.

\subsection{Normalising earnings}

The earnings of a private company reflect its owner's choices as well as its business. Before any
model is fitted, the reported figures are restated to what the company would earn under a buyer who
pays market rates and expects no unusual items. \Cref{tab11:norm} shows the restatement for a stylised
manufacturer of industrial components with revenue of \nmRev{} million. The owner draws pay above what
a hired manager would cost, a legal settlement will not recur, the company rents its plant from the
owner below the market rent, an inventory write-down was a one-off, and a gain on selling equipment
is not operating income. Normalised EBITDA is \nmNorm{} million, \nmChange\% above the reported figure,
a margin of \nmMargin\%. Every one of these adjustments is a judgment, and each should be documented
with the evidence for the market rate it uses; the pay and rent adjustments in particular are ones that
only an owner with control could make, which matters for the discounts of \cref{sec11:discounts}. A buyer should also
separate adjustments any owner could make from those that need control or specific operating changes,
and leave out changes already in the forecast.

\begin{table}[!tb]
  \centering
  \small
  \caption{Normalisation of the stylised manufacturer's EBITDA, in millions of dollars. The company and
  the adjustments are illustrations.}
  \label{tab11:norm}
  \begin{tabular}{@{}lr@{}}
    \toprule
    Reported EBITDA & \nmRep\\
    add the owner's pay above a market salary & \nmPay\\
    add a one-off legal settlement & \nmLegal\\
    less the market rent above the rent paid to the owner & \nmRent\\
    add an inventory write-down not expected to recur & \nmStock\\
    less a gain on the sale of equipment & \nmGain\\
    \midrule
    Normalised EBITDA & \nmNorm\\
    \bottomrule
  \end{tabular}
\end{table}

\subsection{The cost of capital}

The discount rate is built from market inputs. The risk-free rate is the ten year Treasury yield of
September 2026, \wkRf\% \parencite{fredch11}. The equity risk premium is the implied premium of \textcite{damodaran2026}
at the start of 2026, \wkErp\%, computed from the level of the S\&P 500 and its expected cash flows;
\cref{fig11:erp} shows its history since \wkErpFirst, with a mean of \wkErpMean\%, a low of
\wkErpMin\% at the end of \wkErpMinYear{} and a high of \wkErpMax\% at the end of \wkErpMaxYear. That
premium was computed against a Treasury yield of \wkTbond\%, and pairing it with the September yield is
an assumption: the inputs are of different dates, so the cost of capital is a mixed-date illustration
and not one calibrated to a single market day. A private company has no beta of its own, so it borrows one from an industry, here the
\wkFirms{} US machinery companies, with an average levered beta of \wkBetaL{} and debt of \wkDvInd\% of
value. The beta has to be unlevered at the industry's leverage and relevered at the company's target of
\wkTarget\% debt to value, and the two common rules give different answers.

\Textcite{hamada1972} assumes a constant amount of debt, so the tax shields are as safe as the debt:
$\beta_E=\beta_U\bigl(1+(1-\tau)D/E\bigr)$. \Textcite{harris1985} assume that debt is kept at a
constant share of value, so the tax shields have the risk of the assets: $\beta_E=\beta_U(1+D/E)$.
With a marginal tax rate of \wkTax\%, the unlevered beta is \wkBuHam{} by the first rule and \wkBuHp{}
by the second, and the relevered betas are \wkBrHam{} and \wkBrHp. With debt at \wkRd\%, \wkSpread{}
points over the Treasury yield, the usual formula
$W=\frac EV r_E+\frac DV r_D(1-\tau)$ gives \wkWHam\% and \wkWHp\%. The difference is not a matter of
taste. For a company that keeps its leverage at a target, the assumptions of Harris and Pringle are
closer to its financing policy, and the constant debt rule mixes into the beta a policy the company
does not follow; with the target reset once a year rather than continuously, the Miles and Ezzell
treatment below is the more precise one.

\begin{figure}[!tb]
  \centering
  \includegraphics[width=0.84\textwidth]{ch11/fig12_erp}
  \caption{The implied equity risk premium of the S\&P 500 at the end of each year, from
  \textcite{damodaran2026}, and the ten year Treasury yield it is measured against.}
  \label{fig11:erp}
\end{figure}

With a target leverage the right weighted average cost of capital also depends on how often the debt
is reset. \textcite{miles1980} worked it out for a reset once a year.

\begin{proposition}[Miles and Ezzell]\label{prop11:me}
Let a firm's free cash flows have cost of capital $r_U$, and let its debt be reset at the end of each
year, after that year's cash flow, to a share $L$ of its value, at the rate $r_D$, with interest
deductible at $\tau$ and always usable, so that the interest of year $t$ is fixed at $t-1$. Then the value
satisfies $V_{t-1}=(F_t+V_t)/(1+W)$ with
\begin{equation}\label{eq11:me}
  W=r_U-L\,r_D\,\tau\,\frac{1+r_U}{1+r_D}.
\end{equation}
\end{proposition}

\begin{proof}
The tax shield paid at $t$ is $\tau r_DLV_{t-1}$. It is known at $t-1$, so it is discounted for its last
year at $r_D$, while the firm's own cash flow and its value at $t$ have the risk of the assets and are
discounted at $r_U$. The adjusted present value\index{adjusted present value} of \textcite{myers1974} over one year is therefore
$V_{t-1}=(F_t+V_t)/(1+r_U)+\tau r_DLV_{t-1}/(1+r_D)$. Collecting $V_{t-1}$,
$V_{t-1}\bigl(1-\tau r_DL/(1+r_D)\bigr)(1+r_U)=F_t+V_t$, and the factor on the left is $1+W$ with $W$ as
in \cref{eq11:me}. The recursion is a valuation relation under these discount rates; it need not
hold path by path.
\end{proof}

For the manufacturer, the Harris and Pringle unlevered beta gives $r_U=\wkRu\%$, and
\cref{eq11:me} gives $W=\wkWme\%$. The usual formula gives the same number when the cost of equity is
relevered consistently with the annual reset, $r_E=r_U+(r_U-r_D)\frac DE\bigl(1-\tau r_D/(1+r_D)\bigr)$,
which is \wkReMe\%; the code checks the equality. The figure of \wkWHp\% computed above came from a
beta relevered with a debt beta of zero, as if the debt carried no risk premium, while it is charged
\wkRd\%, and with the tax shields given the risk of the assets as under continuous rebalancing. Building the
cost of equity from $r_U$ and $r_D$ directly, with the target reset once a year, removes both
inconsistencies and gives \cref{eq11:me}.

\Cref{fig11:wacc} sets the three treatments of this section against the target leverage. The two textbook rules differ
even with no debt, \wrHamZero\% and \wrHpZero\%, because they unlever the industry beta differently.
With a riskless debt beta, the rule of Harris and Pringle gives $W=r_U+\frac DV\,(r_D-r_f-\tau r_D)$,
which rises slightly with leverage here because the spread of \wkSpread{} points is larger than the tax
saving $\tau r_D$; the rule of Hamada\index{Hamada} falls, and Miles and Ezzell falls fastest. At the \wkTarget\% target
the three are \wrSpreadTarget{} points apart and at 60\% they are \wrSpreadSixty{} points apart, so the
choice matters more the more the company borrows.

\begin{figure}[!tb]
  \centering
  \includegraphics[width=0.8\textwidth]{ch11/fig20_wacc}
  \caption{Cost of capital against target debt to value for the machinery inputs of the text: betas
  unlevered and relevered by Hamada and by Harris and Pringle with the usual formula, and the Miles and
  Ezzell cost for a target reset once a year. The dotted line marks the \wkTarget\% target.}
  \label{fig11:wacc}
\end{figure}

\subsection{The discounted cash flow}

The manufacturer's revenue grows by 5\%, 5\%, 4\%, 4\% and 3\% over five years, at the normalised margin,
with depreciation of 3\% of revenue, capital spending of 3.5\% and working capital of 12\% of the increase
in revenue, and tax at \wkTax\%. The free cash flows are \vlFcfOne, \vlFcfTwo, \vlFcfThree, \vlFcfFour{}
and \vlFcfFive{} million. After year five the company grows at \vlG\% and earns \vlRoc\% on new capital, so
it reinvests \vlReinv\% of NOPAT, and with NOPAT of \vlNopatSix{} million in year six the terminal value
of \cref{eq11:driver} is \vlTv{} million. Discounted at \wkWme\%, the enterprise value is \vlEv{} million,
of which \vlTvShare\% comes from the terminal value, and the equity, after today's net debt of \vlNetDebt{}
million, is worth \vlEq{} million; the cost of capital assumes the target leverage that a buyer would
put in place, whatever the owner borrows today. The enterprise value is \vlEvEbitda{} times normalised EBITDA.

\Cref{tab11:dcf} sets out the forecast line by line, so that the enterprise value can be traced from
revenue to present value.

\begin{table}[!tb]
  \centering
  \small
  \caption{The manufacturer's five year DCF, in millions of dollars, discounted at \wkWme\%. The terminal
  value is that of \cref{eq11:driver} at the end of year five.}
  \label{tab11:dcf}
  \begin{tabular}{@{}lrrrrr@{}}
    \toprule
    Year & 1 & 2 & 3 & 4 & 5\\
    \midrule
    Revenue & \vdRevOne & \vdRevTwo & \vdRevThree & \vdRevFour & \vdRevFive\\
    EBITDA & \vdEbitdaOne & \vdEbitdaTwo & \vdEbitdaThree & \vdEbitdaFour & \vdEbitdaFive\\
    Depreciation & \vdDaOne & \vdDaTwo & \vdDaThree & \vdDaFour & \vdDaFive\\
    EBIT & \vdEbitOne & \vdEbitTwo & \vdEbitThree & \vdEbitFour & \vdEbitFive\\
    NOPAT & \vdNopatOne & \vdNopatTwo & \vdNopatThree & \vdNopatFour & \vdNopatFive\\
    Capital spending & \vdCapexOne & \vdCapexTwo & \vdCapexThree & \vdCapexFour & \vdCapexFive\\
    Working capital & \vdDwcOne & \vdDwcTwo & \vdDwcThree & \vdDwcFour & \vdDwcFive\\
    Free cash flow & \vdFcfOne & \vdFcfTwo & \vdFcfThree & \vdFcfFour & \vdFcfFive\\
    Discount factor & \vdDfOne & \vdDfTwo & \vdDfThree & \vdDfFour & \vdDfFive\\
    Present value & \vdPvOne & \vdPvTwo & \vdPvThree & \vdPvFour & \vdPvFive\\
    \midrule
    \multicolumn{5}{@{}l}{Present value of the five cash flows} & \vdPvSum\\
    \multicolumn{5}{@{}l}{Terminal value \vlTv, discounted} & \vdPvTv\\
    \multicolumn{5}{@{}l}{Enterprise value} & \vlEv\\
    \bottomrule
  \end{tabular}
\end{table}

The same value follows from the adjusted present value: discounting the free cash flows and the
terminal value at $r_U$ gives an unlevered value of \vlVu{} million, and the tax shields of the
rebalanced debt, valued by the recursion in the proof of \cref{prop11:me}, add \vlShield{} million, or
\vlShieldShare\% of the total. The code checks that the two routes agree to rounding error. The shortcut
$\tau D$, which treats a debt of \wkTarget\% of today's value as permanent, would give \vlTauD{} million.
In this illustration that is less than the rebalanced shields are worth, because the debt here grows
with the firm; each year's shield is riskier than a fixed one, which pulls the other way but by less.
The ordering comes from these assumptions about debt growth, shield risk and discounting, and need not
hold elsewhere.

A one stage capitalisation of the first year's free cash flow, $F_1/(W-g)$, gives \vlCapEarn{}
million, more than the DCF. It treats the first year's free cash flow as the numerator, one period ahead, of
a perpetuity growing at \vlG\% from then on. The two disagree because the one stage formula applies the terminal growth to the first year's
cash flow, whose reinvestment is lower than the \vlReinv\% of NOPAT that \cref{prop11:driver} requires
for growth of \vlG\% at a return of \vlRoc\%, and so carries that year's reinvestment and profitability
into every later year. Capitalised earnings are
a valid shortcut only for a company whose current cash flow already reflects the reinvestment its
growth needs.

\Cref{fig11:tornado} shows how the value moves when one input changes, as percentage changes in
enterprise value. A point less and a point more on the cost of capital move it by \vlToWaccUp\% and \vlToWaccDn\%; two points of growth in each forecast year by \vlToNearUp\%
and \vlToNearDn\%; a point of margin by \vlToMarginUp\% and \vlToMarginDn\%; a return on new capital of
15\% or 9\% by \vlToRocUp\% and \vlToRocDn\%; half a point of terminal growth only by \vlToGrowthUp\% and
\vlToGrowthDn\%, because with a return on new capital of \vlRoc\% against a cost of capital of
\wkWme\% growth creates little value (\cref{prop11:driver}); and half a point less or more of capital spending by \vlToCapexUp\% and \vlToCapexDn\%.
A DCF of a private company is sensitive to at least these input changes, and an honest report gives
the range alongside the point estimate.

\begin{figure}[!tb]
  \centering
  \includegraphics[width=0.84\textwidth]{ch11/fig13_tornado}
  \caption{Change in the manufacturer's enterprise value when one input of the DCF is moved, the
  others held at their base values.}
  \label{fig11:tornado}
\end{figure}

\subsection{Multiples}

The market approach applies a multiple paid for comparable companies to the company's own earnings.
The US machinery industry trades at \vlIndMult{} times EBITDA, which on the manufacturer's normalised
EBITDA would give an enterprise value of \vlMktEv{} million, \vlMktGap\% above the DCF. The industry's companies are public,
larger and more diversified than the manufacturer, and they differ from it in liquidity, control and the
date of the data, so a buyer who applies their multiple must either show that growth and risk are
comparable or adjust the multiple for the difference. A multiple is a summary of a DCF that someone else has done, and it is only as
comparable as the forecasts behind it.

\Cref{prop11:driver} turns a multiple into the inputs it implies. With next year's NOPAT in the
denominator, the forward multiple is $\mathit{EV}/\mathit{NOPAT}_{1}=(1-g/\mathit{ROC})/(W-g)$; with this
year's NOPAT, the trailing multiple is $(1+g)(1-g/\mathit{ROC})/(W-g)$, the form used below because the
data report trailing earnings. Either way a multiple fixes a relation between growth, the return on
capital and the cost of capital. For the \mkN{} US industries in the Damodaran data with a return on
capital above 4\% and a cost of capital more than a point above 2.5\%, the multiple of trailing after-tax EBIT implied by each
industry's own cost of capital and return on capital, with growth of 2.5\%, has a median of \mkImpMed,
against an observed median of \mkObsMed; the observed multiple is above the implied one in \mkAbove\%
of industries (\cref{fig11:multiples}). Across industries the two are negatively related, with a rank
correlation of \mkRho, which a five year stage of growth at each industry's own fundamental growth rate
before the terminal value weakens to \mkRhoTwo{} but leaves negative. The observed multiples rise with
the industry's cost of capital, rank correlation \mkRhoW; expected growth, profitability and
differences in composition and measurement may help explain that, though the correlation does not
establish a cause. At the medians, taken one at a time and
unrounded, of a cost of capital of
\mkWmed\% and a return on capital of \mkRocMed\%, the observed trailing multiple implies, under stable-state assumptions, perpetual growth of
\mkGstar\%; the three medians need not belong to any one industry. Reading a market multiple through the value driver is a way of finding the growth that
the price assumes, and judging whether it is plausible for the company being valued.

The R listing below does this for any multiple. It solves the trailing value driver for $g$ with
\texttt{uniroot} and checks the answer against the root of the quadratic that the equation reduces to;
at the medians, typed in to the digits shown, it prints \mkGstar\% both ways.
\Needspace*{16\baselineskip}
\lstinputlisting[language=R]{code/ch11/snippets/implied_growth.R}

\begin{figure}[!tb]
  \centering
  \includegraphics[width=0.78\textwidth]{ch11/fig11_multiples}
  \caption{Observed EV/EBIT$(1-\tau)$ of US industries in January 2026 against the trailing multiple implied by
  \cref{eq11:driver}, $(1+g)(1-g/\mathit{ROC})/(W-g)$, with each industry's cost of capital and return on capital and growth of 2.5\%.
  Data from \textcite{damodaran2026}.}
  \label{fig11:multiples}
\end{figure}

\FloatBarrier
% =====================================================================================
\section{Discounts for control and marketability}
\label{sec11:discounts}

The DCF of \cref{sec11:value} values the whole company as a buyer with control would run it, with the
pay and rent restated. A minority stake in the same company is worth less for two reasons. Its owner
cannot set the pay, the dividends or the date of sale, and there is no market in which to sell it.
The first is the \emph{discount for lack of control}\index{discount for lack of control} (DLOC) and the second the \emph{discount for
lack of marketability\index{marketability}} (DLOM). Because the two discounts multiply, their order does not change the
result; what matters is that the DLOC is taken from a control value and the DLOM from a marketable
minority value, without counting either twice.

If buyers of control pay a premium $p$ over the minority value, the minority value is $1/(1+p)$ of
the control value and $\mathrm{DLOC}=1-1/(1+p)$: a premium of \dmCp\% is a discount of \dmDloc\%. A
premium observed in acquisitions is not automatically the right DLOC for a given minority stake,
because it can include synergies particular to the buyer. The
DLOM is then a discount on the marketable minority value, so the two combine as
$1-(1-\mathrm{DLOC})(1-\mathrm{DLOM})$, which is less than their sum by their product. The size of the DLOM
is the harder question, and two option arguments give benchmarks for it.

\paragraph{The protective put.}
\textcite{chaffe1993} observed that the owner of a share that cannot be sold for $T$ years could
set a minimum on the share's value at the end of that period, while keeping the upside, by buying a
European put at the current price, though
the put does not restore the right to sell at an earlier date. The price of that put, as a share of the
share price, is a model-based proxy for what the restriction costs. It is neither an exact DLOM nor a
discount observed in transactions. In the model of
\textcite{black1973} and \textcite{merton1973} with no dividends, an at the money put on a share of
price one is $e^{-rT}\Phi(-d_2)-\Phi(-d_1)$ with $d_{1,2}=(r\pm\sigma^2/2)\sqrt T/\sigma$. For a volatility of
\dmSig\% and a restriction of \dmT{} years it is \dmChaffe\% at the Treasury yield of \wkRf\%, used as a continuous rate, and
\dmChaffeZero\% at a zero rate. The put cannot cost more than the share, so this measure stays below
100\%, and it treats the holder as if the only loss were the inability to avoid a fall.

\paragraph{The lookback value.}
\textcite{longstaff1995} asked instead what an owner would give up against an investor who could sell
at the best moment during the restriction. Longstaff\index{Longstaff} presented this as an upper bound on the discount, but it is a bound only inside his
frictionless model with perfect timing; read here, it is the value of an idealised right and sets no ceiling
on a DLOM measured in transactions.

\begin{proposition}[Longstaff's lookback value]\label{prop11:longstaff}
Let $Z$ be a standard Brownian motion, let the discounted share price be
$S_t=\exp(\sigma Z_t-\sigma^2t/2)$, and let $v=\sigma^2T$. Then
\begin{equation}\label{eq11:longstaff}
  \E\Bigl[\max_{0\le t\le T}S_t\Bigr]-1=\Bigl(2+\frac v2\Bigr)\Phi\Bigl(\frac{\sqrt v}2\Bigr)
  +\sqrt{\frac v{2\pi}}\,e^{-v/8}-1,
\end{equation}
which depends on $\sigma$ and $T$ only through $v$ and grows without bound in $v$. It is measured against
the initial price and can exceed one.
\end{proposition}

\begin{proof}
Write $X_t=\sigma Z_t-\sigma^2t/2$, a Brownian motion with drift $\nu=-\sigma^2/2$ and variance $\sigma^2$,
and $M=\max_{t\le T}X_t\ge0$. With $s=\sqrt v$, the law of the maximum of a Brownian motion with drift gives
\[
  \mathbb P(M>m)=\Phi\Bigl(\frac{-m+\nu T}{s}\Bigr)+e^{2\nu m/\sigma^2}\Phi\Bigl(\frac{-m-\nu T}{s}\Bigr)
  =\Phi\Bigl(-\frac ms-\frac s2\Bigr)+e^{-m}\Phi\Bigl(-\frac ms+\frac s2\Bigr)
\]
for $m\ge0$. Since $\E[e^M]-1=\int_0^\infty e^m\mathbb P(M>m)\,dm$, there are two integrals. The second term
gives $I_2=\int_0^\infty \Phi(s/2-m/s)\,dm=s\int_{-\infty}^{s/2}\Phi(u)\,du=s\bigl(\tfrac s2\Phi(\tfrac s2)+\varphi(\tfrac s2)\bigr)$.
For the first, integration by parts gives
$I_1=-\Phi(-\tfrac s2)+\tfrac1s\int_0^\infty e^m\varphi(\tfrac ms+\tfrac s2)\,dm$, and since
$e^m\varphi(\tfrac ms+\tfrac s2)=\varphi(\tfrac ms-\tfrac s2)$, the integral is $s\Phi(\tfrac s2)$ and
$I_1=2\Phi(\tfrac s2)-1$. Adding, $I_1+I_2=(2+\tfrac{s^2}2)\Phi(\tfrac s2)+s\varphi(\tfrac s2)-1$, which is
\cref{eq11:longstaff}. As $v\to\infty$ the right side grows like $v/2$.
\end{proof}

For \dmSig\% volatility over \dmT{} years the lookback value is \dmLong\%, \dmLongOverPut{} times the put. The code checks
\cref{eq11:longstaff} by simulating \dmMcPaths{} paths with \dmMcSteps{} steps, drawing the maximum within
each step exactly from its Brownian bridge law; at three horizons and two volatilities the simulated
means lie within \dmMcMaxZ{} standard errors of the formula. \Cref{fig11:dlom} draws both measures. At
30\% volatility the put rises from \dmChThirtyOne\% for a year to \dmChThirtyFour\% for four years and the
lookback value from \dmLsThirtyOne\% to \dmLsThirtyFour\%; at 60\% the lookback value for four years is
\dmLsSixtyFour\%, more than the share itself. The value for 30\% over four years equals that for 60\% over
one, since only $v$ matters. The two models value different idealised rights, so they are not a
calibrated interval for a DLOM: at \dmSig\% over \dmT{} years they give \dmChaffe\% and \dmLong\%. The
choice of measure and of volatility moves the answer more than any other input of the minority valuation below.

\paragraph{An average-strike put.}
Neither benchmark matches a seller with no special timing ability. \textcite{finnerty2012} models the
lost right as a put that pays the average price over the restriction period in exchange for the
price at its end, an investor who sells at random dates during the period receiving on average the
mean price. Approximating the law of the average by a lognormal with matched moments gives
\begin{equation}\label{eq11:finnerty}
  D_F=e^{-\delta T}\bigl(2\Phi(\tfrac{w}{2})-1\bigr),\qquad
  w^2=v+\ln\bigl(2(e^{v}-v-1)\bigr)-2\ln(e^{v}-1),
\end{equation}
with $v=\sigma^2T$ and $\delta$ the dividend yield, here zero. Since $\Phi(w/2)<1$, $D_F$ stays below the
share price whatever the variance, unlike the lookback value. At \dmSig\% over \dmT{} years it gives
\dmFin\%, below the put at a zero rate of \dmChaffeZero\%. A simulation of the average-strike payoff
with 250 steps a path gives \dmFinMc\% (standard error \dmFinMcSE), so the moment-matched closed form
is slightly low; at 60\% over four years the gap widens to \dmFinSixtyFour\% against
\dmFinMcSixFour\%. At 30\% the measure runs from \dmFinThirtyOne\% for a year to \dmFinThirtyFour\% for
four. \textcite{finnerty2012} reports that the model tends to understate the discounts observed in
private placements of restricted stock, although it is a significant explanatory variable for them.
Used in place of the put, it values the minority stake below at \dmStakeFin{} million.

\begin{figure}[!tb]
  \centering
  \includegraphics[width=0.84\textwidth]{ch11/fig14_dlom}
  \caption{Two measures of the discount for lack of marketability against the length of the
  restriction: the price of an at the money European put at the Treasury yield of \wkRf\% and the
  lookback value of \cref{eq11:longstaff}, with simulated values of the lookback.}
  \label{fig11:dlom}
\end{figure}

\paragraph{A minority stake.}
Take a \dmStakeShare\% stake in the manufacturer of \cref{sec11:value}. Its share of the equity at the
control value is \dmStakeCtl{} million. A control premium of \dmCp\% gives a DLOC of \dmDloc\%, and the put
for \dmT{} years at \dmSig\% volatility a DLOM of \dmChaffe\%, so the combined discount is \dmComb\%, not
the \dmSum\% of their sum, and the stake is worth \dmStake{} million. The DLOM here is a
modelling choice taken from the put benchmark, so the value is an illustration and not an observed
fair market value; the control adjustments enter through the DLOC alone. Had the earnings not been normalised
for the owner's pay and rent, part of the DLOC would already have been taken in the cash flows; the
discount should be applied to a control value only if the cash flows are those of a controlling owner.

\Cref{tab11:minority} shows how far the answer moves with the measure, the volatility and the
restriction period, with the DLOC held at \dmDloc\%. With the put, the stake is worth between
\mgPutSixtyFour{} and \mgPutThirtyOne{} million across the grid. With the lookback value it falls to
\mgLbFortyFiveFour{} million at 45\% volatility over four years, and at 60\% over four years the lookback
value is more than the share price, so it cannot serve as a discount at all. The Octave script
\texttt{code/ch11/snippets/minority.m}, kept in the repository and not printed, computes the grid from the
same formulas.

\begin{table}[!tb]
  \centering
  \small
  \caption{Value of the \dmStakeShare\% stake, in millions of dollars, with a DLOC of \dmDloc\% and a DLOM
  from the protective put or from the lookback value. One cell is missing (n.a.): there the lookback value
  is above the share price and gives no usable discount.}
  \label{tab11:minority}
  \begin{tabular}{@{}lrrrrrr@{}}
    \toprule
    & \multicolumn{3}{c}{Protective put} & \multicolumn{3}{c}{Lookback value}\\
    \cmidrule(lr){2-4}\cmidrule(l){5-7}
    Volatility & 1 year & 2 years & 4 years & 1 year & 2 years & 4 years\\
    \midrule
    30\% & \mgPutThirtyOne & \mgPutThirtyTwo & \mgPutThirtyFour & \mgLbThirtyOne & \mgLbThirtyTwo & \mgLbThirtyFour\\
    45\% & \mgPutFortyFiveOne & \mgPutFortyFiveTwo & \mgPutFortyFiveFour & \mgLbFortyFiveOne & \mgLbFortyFiveTwo & \mgLbFortyFiveFour\\
    60\% & \mgPutSixtyOne & \mgPutSixtyTwo & \mgPutSixtyFour & \mgLbSixtyOne & \mgLbSixtyTwo & \mgLbSixtyFour\\
    \bottomrule
  \end{tabular}
\end{table}

\FloatBarrier
% =====================================================================================
\section{Limitations}
\label{sec11:limits}

The stylised fund, the buyout and the manufacturer are illustrations with round inputs, and the
simulations draw deal returns from lognormal laws with one beta and no link between the timing of
exits and the market. Real funds sell more when markets are high, which moves both the J-curve and
the gap between IRR and PME. The NAV of the simulated funds is marked at cost for a year and then at
fair value without smoothing; reported NAVs lag the market, and a PME computed on an interim NAV
inherits the lag.

The CalPERS data cover only the partnerships active at the end of March 2026. A fund that was wound
up before then is missing, so the older vintages are a selection of those that still held assets, and
their multiples may differ from the multiples of every fund of the same vintage. The data give totals
at one date and no timing of the cash flows, so the PME of a CalPERS fund cannot be computed from
them, and the implied duration of \cref{prop11:duration} expresses the reported IRR and TVPI as an
equivalent single period investment, without assuming that the fund had one contribution and one
distribution.

The cost of capital pairs an implied equity premium measured at the start of 2026 with a Treasury
yield of September 2026 and with industry averages of January 2026. The industry data are averages
over companies far larger than the manufacturer, and the market multiples are those of public
companies. The DCF has no size premium and no specific risk premium, which practitioners often add
for small private companies; adding them would lower the model value with everything else unchanged,
and whether they are justified depends on whether those risks are already in the forecasts or the
beta. The DLOM measures are model benchmarks, not estimates fitted to transactions in restricted shares.

\FloatBarrier
% =====================================================================================
\section{Summary}
\label{sec11:summary}

\begin{itemize}
  \item In a whole fund waterfall the LPs get their capital back and a preferred return before the GP
        takes carry; with a full catch-up the GP ends with exactly the carry share of the profit once
        the catch-up is complete, provided no capital is called afterwards. For the stylised fund the gross IRR of \wfGrossIrr\% becomes a net
        IRR of \wfNetIrr\%.
  \item Deal by deal carry is worth more to the GP than whole fund carry with the same terms even when
        the clawback returns every excess dollar, because the GP holds the excess in between. In
        simulated funds the clawback was triggered in \wfTmClawFreq\% of funds.
  \item The J-curve comes from fees and from marking deals at cost while capital is called. Interim
        IRRs rank funds poorly in the first years: the rank correlation with the final IRR is
        \jcRankThree{} after three years and \jcRankEight{} after eight.
  \item The KS-PME divides the index-discounted distributions by the index-discounted contributions,
        and direct alpha is the IRR of the flows carried with the index. For flows that change sign
        once the two agree on whether the fund beat the index. When the market's premium over the risk-free rate
        exceeds its variance, a fund with beta above one shows an expected PME above one with no skill;
        with the simulation's index but leverage financed at the Treasury yield of \wkRf\%, the premium is
        smaller than the variance and the bias runs the other way, an expected PME of \pkBetaTwo{} at a beta of two over \pkT{} years.
  \item In a buyout the equity multiple is $(U-\lambda \kappa)/(1-\lambda)$, so debt helps exactly when the
        asset multiple beats the compounded after-tax cost of debt. The stylised deal returned \lxIrr\%
        with debt against \lxIrrU\% without.
  \item The value driver formula ties terminal growth to reinvestment: growth adds value only when
        the return on new capital exceeds the cost of capital.
  \item With debt reset once a year to a target share of value, the Miles and Ezzell cost of capital
        and the adjusted present value give the same value; the manufacturer is worth \vlEv{} million as
        a whole, with \vlTvShare\% in the terminal value.
  \item Discounts for control and marketability combine multiplicatively. The protective put and
        the lookback value of \cref{prop11:longstaff} give \dmChaffe\% and \dmLong\% for a restriction of
        \dmT{} years at \dmSig\% volatility, outputs of two different models and not lower and upper
        estimates of the discount.
\end{itemize}

\FloatBarrier
% =====================================================================================
\section{Exercises}
\label{sec11:exercises}

\begin{enumerate}[label=\textbf{11.\arabic*}, ref=11.\arabic*, leftmargin=*]
  \item\label{exr11:one} A fund calls 100 at time zero and distributes 60 after three years and 120 after
  six. The terms are a preferred return of 8\% compounded annually, carry of 20\% and a full catch-up, in a
  whole fund waterfall. Split each distribution into return of capital, preferred return, catch-up and
  the 80/20 split. What do the GP and the LPs receive in total, and what is the LPs' IRR?
  \item\label{exr11:two} Under deal by deal carry a GP takes carry of 10 on a deal sold after two years.
  When the fund ends after five years the whole fund carry is 4, and the clawback returns the rest. With
  a continuously compounded rate of 8\%, what is the present value to the GP of the deal by deal terms
  over the whole fund terms, if the clawback is paid in full?
  \item\label{exr11:three} A fund reports a TVPI of 1.8 and a net IRR of 15\%. What is its implied
  duration?
  \item\label{exr11:four} A fund calls 100 when an index stands at 100 and returns 160 four years later,
  when the index stands at 130. Find the KS-PME, the direct alpha and the IRR.
  \item\label{exr11:five} The index has a continuous drift of 7\% and a volatility of 16\%, and the
  continuously compounded risk-free rate is 4\%. A fund with no alpha calls one dollar and holds a position of $\beta$ in the
  index, rebalanced continuously and financed at the risk-free rate, for five years. What is its
  expected PME at $\beta=1.5$ and at $\beta=0.5$?
  \item\label{exr11:six} Debt costs 7\% before a tax of 25\% and a buyout lasts five years. Find $\kappa$. With
  debt of 60\% of the price, what is the equity multiple if the asset multiple is 1.8, and if it is 1.2?
  \item\label{exr11:seven} A company is bought at 9 times EBITDA of 40 with net debt of 200 and sold at 10
  times EBITDA of 55 with net debt of 120. Split the equity gain into growth, the multiple and debt
  repaid with \cref{lem11:bridge}.
  \item\label{exr11:eight} Next year's NOPAT is 10, the cost of capital 8\% and the return on new capital
  15\%. Find the value at growth of 3\% and of 5\%, and without growth.
  \item\label{exr11:nine} Assets have a cost of capital of 9\%, debt costs 6\%, the tax rate is 25\%, and
  debt is kept at 40\% of value. Compare the Miles and Ezzell cost of capital with
  $r_U-L\,r_D\,\tau$, the rate for continuous rebalancing.
  \item\label{exr11:ten} A comparable company has a levered beta of 1.3 and a ratio of debt to equity of
  0.5, with a tax rate of 25\%. Find its unlevered beta under the rules of Hamada and of Harris and
  Pringle.
  \item\label{exr11:eleven} A private company reports EBITDA of 8.0 million. Its owner draws 0.9 million,
  where a hired manager would cost 0.35 million, and the year includes a one-off gain of 0.3 million.
  Find normalised EBITDA and the value at 7 times.
  \item\label{exr11:twelve} The control premium is 30\% and the DLOM 25\%. Find the DLOC, the combined
  discount and the value of a 15\% stake in a company whose equity is worth 60 million under control.
  \item\label{exr11:thirteen} For a restriction of two years on a share with a volatility of 40\%, find
  the protective put at a continuously compounded rate of 4\% and the lookback value of \cref{prop11:longstaff}.
  \item\label{exr11:fourteen} A business earns 7\% on new capital and has a cost of capital of 9\%. Find
  the value, as a multiple of next year's NOPAT, at growth of 2\% and 4\%, and without growth. Why does faster growth
  lower the value?
\end{enumerate}

\subsection*{Computing projects}
\label{sec11:projects}
These use the data and code in the repository. Each is solved by a function in
\texttt{code/ch11/projects.py}, which \texttt{test\_projects.py} runs, so an answer can be
checked with \texttt{python code/ch11/projects.py} from the root of the repository. The
shipped series end in September 2026, and those that end earlier are noted.

\begin{enumerate}[label=\textbf{P11.\arabic*}, leftmargin=*, ref=P11.\arabic*]
\item Rerun the buyout of \cref{sec11:intcap} for caps on deductible interest of 10\%, 20\%, 30\% and 40\% of
EBITDA, with and without a carryforward, and debt of four, five and six times EBITDA. Report the loss of
IRR in basis points. Where does the carryforward matter, and at what leverage does a 30\% cap start to
cost more than a few basis points?
\item Tabulate the protective put at a zero rate, Finnerty's average-strike put and Longstaff's lookback value
for volatilities of 20\% to 80\% and restrictions of one to five years with \texttt{chaffe\_dlom},
\texttt{finnerty\_dlom} and \texttt{longstaff\_bound}. In which cells does the lookback value exceed the
share price, and how wide is the range of discounts for a given cell?
\end{enumerate}

\subsection*{Answers}
\begin{enumerate}[label=\textbf{11.\arabic*}, leftmargin=*]
  \item The first distribution is all return of capital, \yaRocThree, and leaves a hurdle balance of
  $100\times1.08^3-60=\yaBalThree$, which grows to \yaBalSix{} by year six. The second returns the remaining
  \yaRocSix{} of capital, pays \yaPref{} of preferred return and \yaCatch{} of catch-up, which solves
  $y=0.2(\yaPref+y)$, and splits the remaining $120-\yaRocSix-\yaPref-\yaCatch=\yaRest$ 80/20. The GP receives \yaGp, 20\% of the profit of 80, the LPs
  \yaLp, and their IRR is \yaLpIrr\%.
  \item The deal by deal terms pay 10 at year two and claw back 6 at year five, while the whole fund terms
  pay 4 at year five. By \cref{prop11:timing}
  the present value advantage is $10(e^{-0.16}-e^{-0.40})=\ybPv$, since the deal by deal path is worth
  $10e^{-0.16}-6e^{-0.40}$ and the whole fund path $4e^{-0.40}$.
  \item $\ln1.8/\ln1.15=\ycDur$ years.
  \item The KS-PME is $(160/130)/(100/100)=\ydPme$. Carried to year four, the call is 130 and the
  distribution 160, so the direct alpha is $\ln(160/130)/4=\ydAlpha\%$ a year, continuously compounded. The
  IRR is $1.6^{1/4}-1=\ydIrr\%$.
  \item By \cref{eq11:betapme}, $\exp\bigl((\beta-1)(0.07-0.04-0.0256)\times5\bigr)$, which is \yePme{} at
  $\beta=1.5$ and \yePmeLow{} at $\beta=0.5$.
  \item $\kappa=(1+0.07\times0.75)^5=\yfK$. By \cref{eq11:lbo} the equity multiple is
  $(1.8-0.6\times\yfK)/0.4=\yfMoic$, and at an asset multiple of 1.2 it is \yfMoicLow, below the asset
  multiple because $1.2<\kappa$.
  \item Growth is $9.5\times15=\ygG$, the multiple $47.5\times1=\ygM$ and debt repaid \ygD, which add to the
  gain of $430-160=\ygTot$.
  \item By \cref{eq11:driver}, $10(1-0.03/0.15)/0.05=\yhV$ at 3\%, \yhVhigh{} at 5\% and $10/0.08=\yhVflat$
  without growth.
  \item By \cref{eq11:me}, $W=0.09-0.40\times0.06\times0.25\times1.09/1.06=\yiMe\%$, against \yiHp\% for continuous
  rebalancing. Each year's shield is fixed a year ahead and discounted at $r_D$ for its last year, so it is a
  little safer than under continuous rebalancing, worth a little more, and the cost of capital is lower.
  \item Hamada gives $1.3/(1+0.75\times0.5)=\yjHam$ and Harris and Pringle $1.3/1.5=\yjHp$.
  \item $8.0+0.55-0.3=\ykE$ million, and the value is \ykV{} million.
  \item The DLOC is $1-1/1.3=\ylDloc\%$, the combined discount $1-(1-0.231)(0.75)=\ylComb\%$, and the stake
  is worth $0.15\times60\times(1-0.423)=\ylStake$ million.
  \item The put is \ymPut\% of the share price, and the lookback value, with $v=0.32$, is \ymLs\%.
  \item By \cref{eq11:driver} the multiples of NOPAT are \ynLow{} at 2\% and \ynHigh{} at 4\%, against
  $1/0.09=\ynFlat$ without growth. Each dollar reinvested earns 7\% and is worth $0.07/0.09$ of a dollar,
  so the more the business grows, the more value it destroys (\cref{prop11:driver}).
\end{enumerate}

\FloatBarrier
% =====================================================================================
\section{Notes, sources and credits}
\label{sec11:credits}

\paragraph{Sources}
The fund terms and the economics of fees and carry are set out in \textcite{metrick2010}, and the
cash flow model of a fund behind the J-curve is in the spirit of \textcite{takahashi2002}. The public
market equivalent is from \textcite{kaplan2005}, direct alpha from \textcite{gredil2023}, the effect of
beta on the PME from \textcite{sorensen2015}, and its generalisation to a stochastic discount factor
from \textcite{korteweg2016}. The evidence on fees and on fund performance is in
\textcite{phalippou2009} and \textcite{harris2014}, the pitfalls of the IRR in \textcite{phalippou2008}, the operating changes after
buyouts in \textcite{kaplan1989}, and the link of leverage to credit markets in \textcite{axelson2013}.
The irrelevance of leverage and the value of the tax shield are from \textcite{modigliani1958,modigliani1963},
the adjusted present value from \textcite{myers1974}, and the cost of capital with annual rebalancing from
\textcite{miles1980}. The rules for unlevering betas are from \textcite{hamada1972} and \textcite{harris1985},
the growing perpetuity from \textcite{gordon1956}, the protective put from \textcite{chaffe1993}, priced
with the model of \textcite{black1973} and \textcite{merton1973}, the average-strike put from
\textcite{finnerty2012}, and the lookback value from \textcite{longstaff1995}. The fixed ratio rule for
interest is described in \textcite{oecd2015}, and the Canadian ratio is in section 18.2 of the
Income Tax Act \parencite{ita182}. The waterfall share and timing results, the duration bound, the sign agreement
of the PME and direct alpha, the leverage identity, the bridge, the value driver and the proof of the
lookback formula are short derivations given here for completeness.

\paragraph{Software used}
The scripts call NumPy \parencite{numpy} for arrays and random numbers, SciPy \parencite{scipy} for root
finding and the normal law, pandas \parencite{pandas} to read the tables, xlrd \parencite{xlrd} to read the
spreadsheets of the Damodaran data, and Matplotlib \parencite{matplotlib} for the figures. The
cross-checks run in R \parencite{rcore} and GNU Octave \parencite{octave}.
\texttt{code/requirements-lock.txt} records the full environment.

\paragraph{Projects for comparison}
The numpy-financial package \parencite{numpyfinancial} and pyxirr \parencite{pyxirr} compute the IRRs of
irregular cash flows used throughout the chapter, and are a check on the code here. The full list, with
links, is in \texttt{CREDITS.md}.

\paragraph{Data}
The partnership data are from the Private Equity Program Fund Performance Review of CalPERS, as of 31
March 2026 \parencite{calpers2026}. The industry averages, the costs of capital and the implied equity
risk premium are from the data sets of Aswath Damodaran, in the update of January 2026
\parencite{damodaran2026}. Both were retrieved on 19 October 2026. They are free to read but not in the
public domain, so they are not redistributed, and \texttt{scripts/pull\_ch11.py} downloads them and
rebuilds the chapter's data files. The ten year Treasury yield is from the Federal Reserve, as published
on FRED \parencite{fredch11}, and ships with the code. The fund, the buyout, the manufacturer and the
exercise inputs are illustrations.

\end{chapterrefs}
\flushbottom
